% !TeX program = xelatex
% !TeX encoding = UTF-8
% =====================================================================
%  成都市初一英语书面表达范文集
%  人教新版（2024 年版）七年级下册  Unit 1 — Unit 8
%  ---------------------------------------------------------------------
%  编译方式（必须使用 XeLaTeX，并连续编译两遍以生成目录与页码）：
%      xelatex 成都市初一英语Unit1-8范文集.tex
%      xelatex 成都市初一英语Unit1-8范文集.tex
%  若使用 IDE（TeXstudio / VS Code + LaTeX Workshop），将编译器设为 XeLaTeX 即可。
% =====================================================================
\documentclass[UTF8,a4paper,11pt,oneside]{ctexbook}

% ---------------- 页面与字体 ----------------
\usepackage{geometry}
\geometry{top=2.5cm,bottom=2.5cm,left=2.4cm,right=2.2cm,headheight=17pt}
\usepackage{xcolor}
\definecolor{mainblue}{RGB}{21,73,124}
\definecolor{accent}{RGB}{176,84,32}
\definecolor{textgray}{RGB}{96,101,107}
\definecolor{rulelight}{RGB}{168,193,219}
\definecolor{patternbg}{RGB}{245,249,253}
\definecolor{transbg}{RGB}{247,247,245}
\definecolor{transframe}{RGB}{176,176,168}

% ---------------- 版式组件 ----------------
\usepackage{fancyhdr}
\usepackage{enumitem}
\usepackage{booktabs}
\usepackage{tabularx}
\usepackage{array}
\usepackage{needspace}
\usepackage[most]{tcolorbox}
\usepackage[hidelinks,bookmarksopen]{hyperref}
\hypersetup{
  pdftitle={Chengdu Grade 7 English Writing Model Essays (Unit 1-8)},
  pdfsubject={中考书面表达备考},
  pdfcreator={XeLaTeX}
}

% ---------------- 章节标题样式 ----------------
\ctexset{
  chapter = {
    format    = {\centering\sffamily\Huge\bfseries\color{mainblue}},
    beforeskip = 1.6ex,
    afterskip  = 2.2ex
  }
}

% ---------------- 页眉页脚 ----------------
\pagestyle{fancy}
\fancyhf{}
\fancyhead[L]{\small\sffamily\color{textgray}\leftmark}
\fancyhead[R]{\small\sffamily\color{mainblue}\thepage}
\fancyfoot[C]{\footnotesize\sffamily\color{textgray}成都市初一英语书面表达范文集（人教新版七年级下册）}
\renewcommand{\headrulewidth}{0.4pt}
\renewcommand{\headrule}{%
  \hbox to\headwidth{\color{rulelight}\leaders\hrule height \headrulewidth\hfill}}
% plain 页面样式：直接用底层定义，避免 fancyhdr v5 的 \fancypagestyle
% 在每页重置全局页眉页脚设置（会把 fancy 的设置清掉）。
\makeatletter
\def\ps@plain{%
  \let\@mkboth\@gobbletwo
  \let\@oddhead\@empty
  \let\@evenhead\@empty
  \def\@oddfoot{\hfil{\small\sffamily\color{textgray}\thepage}\hfil}%
  \let\@evenfoot\@oddfoot
}
\makeatother

% ---------------- 自定义盒子 ----------------
\newtcolorbox{analysisbox}{
  enhanced,
  colback=patternbg, colframe=mainblue!45,
  boxrule=0.5pt, arc=1.2mm,
  left=2.4mm, right=2.4mm, top=1.2mm, bottom=1.2mm,
  toptitle=0.9mm, bottomtitle=0.9mm,
  coltitle=white, colbacktitle=mainblue!78,
  fonttitle=\sffamily\bfseries\small,
  title={亮点解析},
  before skip=1.2ex, after skip=1.6ex,
  fontupper=\small
}
\newtcolorbox{transbox}{
  enhanced,
  colback=transbg, colframe=transframe!75,
  boxrule=0.5pt, arc=1.2mm,
  left=2.4mm, right=2.4mm, top=1.2mm, bottom=1.2mm,
  toptitle=0.9mm, bottomtitle=0.9mm,
  coltitle=white, colbacktitle=transframe!85,
  fonttitle=\sffamily\bfseries\small,
  title={参考译文},
  before skip=1.2ex, after skip=1.6ex,
  fontupper=\small
}

% ---------------- 单元 / 小节 / 范文标题 ----------------
% \fanunit{序号}{英文单元名}{中文单元名}{核心语法}
\newcommand{\fanunit}[4]{%
  \chapter*{\texorpdfstring{Unit #1\quad #2}{Unit #1 #2}}%
  \phantomsection
  \addcontentsline{toc}{chapter}{\texorpdfstring{Unit #1\quad #2（#3）}{Unit #1 #2}}%
  \markboth{Unit #1\quad #2}{Unit #1\quad #2}%
  \vspace{-1.1em}%
  {\centering{\sffamily\large\color{accent}#3}\par}%
  \vspace{0.45em}%
  {\centering\color{rulelight}\rule{0.82\linewidth}{1pt}\par}%
  \vspace{0.75em}%
  {\centering{\small\sffamily\color{textgray}本单元语法重点：#4}\par}%
  \addvspace{1.2em}%
}
% \fansection{〇、本单元写作常用句型} 等
\newcommand{\fansection}[1]{%
  \par\addvspace{2.4ex}%
  \Needspace{6\baselineskip}%
  \noindent{\sffamily\Large\bfseries\color{mainblue}#1}\par
  \nopagebreak
  \addvspace{0.2ex}%
  \noindent{\color{rulelight}\rule{\linewidth}{0.9pt}}\par
  \addvspace{0.9ex}%
  \nobreak
}
% \fanessay{范文一}{中文标题}{English Title}{约126词}
\newcommand{\fanessay}[4]{%
  \Needspace{7\baselineskip}%
  \par\addvspace{2.2ex}%
  \noindent{\color{rulelight}\rule{\linewidth}{0.6pt}}\par
  \nopagebreak
  \addvspace{0.45ex}%
  \noindent{\sffamily\bfseries\color{accent}\large #1}\hspace{0.85em}%
  {\sffamily\bfseries\large\color{mainblue}#2}\par
  \nopagebreak
  \addvspace{0.3ex}%
  {\centering{\bfseries #3}\hspace{0.7em}{\small\color{textgray}（#4）}\par}%
  \addvspace{0.7ex}%
  \nobreak
}
% \fanpat{英文句}{中文释义}{语法说明}
\newcommand{\fanpat}[3]{%
  \item{\bfseries #1}\par\nopagebreak
  {\small\color{textgray}#2}\par\nopagebreak
  {\footnotesize\color{accent!90!black}【用法】#3}\par
}

% ---------------- 写作常用句型列表 ----------------
\newlist{fanpatterns}{enumerate}{1}
\setlist[fanpatterns]{%
  label={\sffamily\bfseries\color{mainblue}\arabic*.},
  leftmargin=2.6em, labelsep=0.7em,
  itemsep=1.0ex, topsep=0.5ex, parsep=0.3ex
}


\begin{document}
\linespread{1.32}\selectfont

\begin{titlepage}
  \thispagestyle{empty}
  \centering
  \vspace*{1.4cm}
  {\color{mainblue}\rule{\linewidth}{1.8pt}}\par
  \vspace{0.55cm}
  {\color{rulelight}\rule{\linewidth}{0.6pt}}\par
  \vspace{1.5cm}
  {\sffamily\bfseries\fontsize{36pt}{46pt}\selectfont\color{mainblue} 七年级下册英语\par}
  \vspace{0.35cm}
  {\sffamily\bfseries\fontsize{36pt}{46pt}\selectfont\color{mainblue} 书面表达范文集\par}
  \vspace{1.3cm}
  {\sffamily\fontsize{15pt}{21pt}\selectfont\color{accent} 人教版（2024 版）· 七年级下册\par}
  \vspace{0.3cm}
  {\sffamily\Large\color{mainblue} Unit 1 — Unit 8（全 8 个单元）\par}
  \vspace{1.0cm}
  {\color{accent}\rule{0.62\linewidth}{0.8pt}}\par
  \vspace{0.9cm}
  {\sffamily\large\color{textgray} 8 个单元\quad·\quad 80 篇范文\quad·\quad 64 条写作常用句型\par}
  \vfill
  {\large\sffamily 适用对象：七年级学生\quad 书面表达备考\par}
  \vspace{0.9cm}
  {\color{rulelight}\rule{0.62\linewidth}{0.6pt}}\par
  \vspace{0.5cm}
  {\sffamily\normalsize\color{textgray} 2026 年 10 月 20 日\par}
  \vspace*{0.6cm}
\end{titlepage}

\frontmatter
\pagestyle{plain}

\chapter*{使用说明}
\phantomsection
\addcontentsline{toc}{chapter}{使用说明}

\noindent\textbf{一、教材依据}\quad 本书依据人教新版（2024 年版）七年级下册编写，覆盖全册
Unit 1 — Unit 8 共 8 个单元，供成都市初一学生中考英语书面表达备考使用。

\vspace{0.6em}
\noindent\textbf{二、单元结构}\quad 每个单元由三部分组成：
\begin{itemize}[leftmargin=2.2em,itemsep=0.25ex,topsep=0.35ex]
  \item \textbf{〇、本单元写作常用句型}：精选 8 条紧扣单元语法的高分句型，附中文释义与用法说明；
  \item \textbf{一、满分范文}：A 组 5 篇，每篇约 120 — 150 词，词汇丰富、句式多样，配\textbf{亮点解析}与\textbf{亮点句}，供冲刺满分的同学精读仿写；
  \item \textbf{二、高分好背好文}：B 组 5 篇，每篇约 85 — 105 词，句式简单、模板化，配\textbf{参考译文}，基础薄弱的同学可直接背诵套用。
\end{itemize}

\vspace{0.6em}
\noindent\textbf{三、使用建议}
\begin{enumerate}[leftmargin=2.2em,itemsep=0.25ex,topsep=0.35ex]
  \item 先背熟本单元的 8 条常用句型与 B 组范文，做到“写得对、写得全”；
  \item 再精读 A 组范文，把「亮点解析」中的加分句式（非谓语动词、定语从句、倒装、\textit{so...that...} 等）迁移到自己的作文里；
  \item 每篇范文至少仿写一次；写完后对照「亮点解析」自查是否用上了本单元语法重点；
  \item 范文话题贴近成都生活（天府绿道、成都大运会、城市音乐厅等），考场上可根据题目灵活替换专有名词。
\end{enumerate}

\clearpage
\vspace{0.2em}
\noindent\textbf{四、单元一览}

\begin{center}
\small
\renewcommand{\arraystretch}{1.25}
\begin{tabularx}{\textwidth}{@{}>{\raggedright\arraybackslash}p{4.6cm}
                              >{\raggedright\arraybackslash}p{6.3cm}
                              >{\raggedright\arraybackslash}X@{}}
\toprule
{\sffamily\bfseries 单元} & {\sffamily\bfseries 核心语法} & {\sffamily\bfseries 代表写作任务}\\
\midrule
{\bfseries Unit 1}\newline Animal Friends\newline {\small 动物朋友} & 一般现在时；can 表能力；形容词描述 & My Animal Friend（我的动物朋友） \\
{\bfseries Unit 2}\newline No Rules, No Order\newline {\small 没有规则就没有秩序} & 祈使句；情态动词 must/have to/can；规则表达 & My Class Rules（我们的班规） \\
{\bfseries Unit 3}\newline Keep Fit\newline {\small 保持健康} & 一般现在时；频率副词；should/can & How I Keep Fit（我如何保持健康） \\
{\bfseries Unit 4}\newline Eat Well\newline {\small 吃得健康} & 可数与不可数名词；there be；量词 & Healthy Eating（健康饮食） \\
{\bfseries Unit 5}\newline Here and Now\newline {\small 此时此刻} & 现在进行时（be doing）；描述正在发生的场景 & What Are They Doing?（此刻正在发生的事） \\
{\bfseries Unit 6}\newline Rain or Shine\newline {\small 无论晴雨} & 天气表达；if 条件句；一般将来时 & The Weather and My Day（天气与我的日子） \\
{\bfseries Unit 7}\newline A Day to Remember\newline {\small 难忘的一天} & 一般过去时；时间顺序与叙事 & A Day to Remember（难忘的一天） \\
{\bfseries Unit 8}\newline Once upon a Time\newline {\small 很久以前} & 一般过去时；故事叙述；once upon a time & A Story I Like（我喜欢的一个故事） \\
\bottomrule
\end{tabularx}
\end{center}

\vspace{0.4em}
\noindent{\small\color{textgray}说明：每单元 A 组 5 篇 + B 组 5 篇，共 10 篇范文；全书共 8 个单元、80 篇范文、64 条写作常用句型。}

\clearpage
\tableofcontents
\mainmatter
\pagestyle{fancy}

% ============================ 正文 ============================

\fanunit{1}{Animal Friends}{动物朋友}{一般现在时；can 表能力；形容词描述}

{\centering\small\color{accent}本单元写作任务：My Animal Friend（我的动物朋友）\par}
\vspace{0.6em}

\fansection{〇、本单元写作常用句型}
\begin{fanpatterns}
\fanpat{My favourite animal is the panda, because it is cute, quiet and a little shy.}{我最喜欢的动物是熊猫，因为它可爱、安静，还有点害羞。}{一般现在时点明喜好，because 引导原因状语从句，三个形容词（短语）并列描写性格，最适合放在开头第一句。}

\fanpat{It has two big black eyes and a short tail, and it looks like a soft ball of fur.}{它有两只大大的黑眼睛和一条短尾巴，看起来像一个毛茸茸的软球。}{has 描写外形，look like 表示“看起来像”；先把眼睛、尾巴写清楚，再打一个比方，动物的样子立刻活起来。}

\fanpat{My dog can understand me when I am sad, and he can make me laugh again.}{我难过的时候，我的狗能懂我，还能让我重新笑起来。}{can + 动词原形表示能力，when 引导时间状语从句，两个 can 句连用，写出动物与人之间的感情。}

\fanpat{Every morning, I take him for a walk along the greenway near my home.}{每天早上，我带着它在我家附近的绿道上散步。}{一般现在时配合 Every morning 表示习惯；take sb for a walk 是“带某人散步”的固定搭配，写宠物日常最好用。}

\fanpat{The monkeys are really funny, and they jump from tree to tree to ask us for food.}{猴子们真有趣，它们在树间跳来跳去，还向我们要吃的。}{先用形容词 funny 概括印象，再用 and 补充具体动作，动词不定式 to ask 表示目的，是动物园见闻类段落的常用写法。}

\fanpat{Animals are our friends, and we can do a lot to help them.}{动物是我们的朋友，我们能做很多事来帮助它们。}{结尾点题句，can 表示能力与可能，to help 作目的状语，适合倡议类、演讲类作文的收束。}

\fanpat{If we keep the water clean, the fish in the river can live a happy life.}{如果我们保持水质干净，河里的鱼就能过上快乐的生活。}{if 条件句用一般现在时，主句用 can 表示结果，写“保护动物”的建议时既讲道理又提做法。}

\fanpat{Not only do animals make us happy, but they also teach us to be kind.}{动物不仅让我们快乐，还教会我们善良。}{Not only 置于句首时主句要部分倒装，but also 补充递进，放在结尾能明显提分。}

\end{fanpatterns}

\fansection{一、满分范文}

\fanessay{范文一}{小狗豆豆}{My Dog Doudou}{约112词}
My favourite animal is my dog Doudou. He is a small brown dog with big black eyes, and he is the newest member of my family.\par
Doudou is only one year old, but he is very clever. He can sit down, shake hands and catch a ball. Every morning, I take him for a walk along the greenway near my home. He runs after birds happily, but he never hurts them. When I come home from school, he waits for me at the door and wags his tail.\par
Doudou is more than a pet to me. He is my best friend, and he makes every day of my life warm and interesting.\par
\begin{analysisbox}
1）开头用 My favourite animal is my dog Doudou 直接点题，with big black eyes 一句勾出小狗的外形，让人一眼记住它。2）主体先用 but he is very clever 一转，再以一个 can 句连写坐下、握手、接球三个本领；Every morning 与 When I come home from school 两个时间状语把日常写得有条理，never hurts them 又写出小狗的温顺。3）结尾用 more than a pet 重新定义人狗关系，以 warm and interesting 收束，感情真挚。
\par\smallskip{\sffamily\bfseries\color{mainblue}亮点句：}{\itshape When I come home from school, he waits for me at the door and wags his tail.}\quad{\small\color{textgray}（译文：我放学回家时，它就在门口等我，摇着尾巴。）}
\end{analysisbox}

\fanessay{范文二}{动物园里的朋友们}{Friends at the Zoo}{约123词}
Chengdu Zoo is a green and quiet place in the north of our city. I often go there with my parents on sunny weekends.\par
We always start with the panda house. The pandas are lazy and lovely, and they eat bamboo slowly and then sleep in the trees. Next to them live the monkeys. They are really funny, and they jump from tree to tree to ask us for food. The elephants are the biggest animals there. They walk slowly and use their long noses to pick up food. My little brother loves the birds best, because their colours are so bright.\par
A day at the zoo is never boring. The animals are our friends, and they teach us to love the world.\par
\begin{analysisbox}
1）开头交代动物园的位置和常去的时间，green and quiet 两个形容词让画面一下子清晰起来。2）主体按“熊猫馆—猴子—大象—鸟”的顺序展开，lazy and lovely、really funny、the biggest 各有侧重；Next to them live the monkeys 用倒装句避免平铺直叙，句式有了变化。3）结尾用 never boring 总结一天感受，再以动物教会我们爱世界收束，主题自然升华。
\par\smallskip{\sffamily\bfseries\color{mainblue}亮点句：}{\itshape Next to them live the monkeys.}\quad{\small\color{textgray}（译文：紧挨着它们住的是猴子。）}
\end{analysisbox}

\fanessay{范文三}{请爱护我们的动物朋友}{Let's Protect Our Animal Friends}{约123词}
Dear friends, do you know that many animals are in danger now? Some of them lose their homes because people cut down too many trees.\par
We can do a lot to help them. First, we can keep the rivers and parks clean. A clean river is a safe home for fish and birds. Second, we can be kind to the animals around us. Please don't throw stones at street cats or birds, because they also feel pain. Third, we can tell our family and friends to say no to animal shows. Animals are our friends, not our toys.\par
Small steps make a big difference. If everyone gives a little love, our city will be a warmer home for every animal. Let's start today.\par
\begin{analysisbox}
1）开头以 Dear friends 呼告，再用问句 do you know that many animals are in danger now 引起读者注意，because 从句交代动物失去家园的原因。2）主体用 First、Second、Third 三条建议层层推进，每条建议后都跟一句理由或做法；A clean river is a safe home for fish and birds 把干净的河流比作安全的家，Please don't throw stones 用祈使句提出行为要求。3）结尾用 If everyone gives a little love 展望美好结果，再以 Let's start today 发出号召，简短有力。
\par\smallskip{\sffamily\bfseries\color{mainblue}亮点句：}{\itshape Animals are our friends, not our toys.}\quad{\small\color{textgray}（译文：动物是我们的朋友，不是我们的玩具。）}
\end{analysisbox}

\fanessay{范文四}{车底下的那只小猫}{The Little Cat Under the Car}{约130词}
It is a cold Saturday morning in December. My cousin and I walk to the park near our school. Suddenly, we hear a weak cry under a parked car.\par
I look down and find a little cat under the car. She is small and wet, and her eyes are full of fear. She cannot stand up. I take off my scarf and wrap her in it. My cousin buys some cat food from a shop, and we feed her slowly. Then we take her to an animal hospital. The doctor says she is only two months old.\par
Now she lives with us, and we call her Lucky. She is healthy and happy, and she follows me everywhere. I am lucky too, because she brings so much joy to my family.\par
\begin{analysisbox}
1）开头用 It is a cold Saturday morning in December 交代场景，Suddenly 引出一声微弱的叫声，一下把读者带进故事。2）主体全部用一般现在时叙述，look down、take off、wrap、feed 几个动作一气呵成，节奏紧凑；small and wet 与 full of fear 写出小猫的无助，为下文救助作铺垫。3）结尾用 I am lucky too 一转，把“我救了小猫”写成“我也被小猫治愈”，立意温暖。
\par\smallskip{\sffamily\bfseries\color{mainblue}亮点句：}{\itshape I take off my scarf and wrap her in it.}\quad{\small\color{textgray}（译文：我解下围巾，把她裹了起来。）}
\end{analysisbox}

\fanessay{范文五}{动物是我们的朋友}{Animals Are Our Friends}{约117词}
Good morning, everyone! Today I want to talk about our animal friends.\par
Animals are around us every day. Some live in our homes, like dogs and cats. Some live in the forests and rivers. They are different in size and colour, but they are all important. Dogs can help blind people walk safely. Birds can tell us when the seasons change. Even small bees work hard and give us sweet honey. Animals are not only beautiful but also helpful to us.\par
However, some people still hurt animals for money. It is time to stop this. Let's be kind to every animal around us. A safe world for animals is a better world for us all. Thank you!\par
\begin{analysisbox}
1）演讲稿开头用 Good morning, everyone! 打招呼，再用 Today I want to talk about... 点明话题，符合演讲的文体要求。2）主体先用 Some... Some... 分写家养动物和野生动物，再用 Dogs can... Birds can... 等 can 句举例说明动物对人的帮助，not only... but also 一句小结，层次清楚。3）结尾用 However 转折指出伤害动物的现象，最后以 A safe world for animals is a better world for us all 收束，把爱护动物与爱护人类自己联系起来，立意更深。
\par\smallskip{\sffamily\bfseries\color{mainblue}亮点句：}{\itshape A safe world for animals is a better world for us all.}\quad{\small\color{textgray}（译文：一个让动物安全的世界，对我们所有人来说也是更好的世界。）}
\end{analysisbox}

\fansection{二、高分好背好文（适合基础薄弱同学）}

\fanessay{范文六}{我家的新成员}{A New Member of My Family}{约91词}
There is a new member in my family now. It is a white rabbit, and we call it Tuantuan.\par
Tuantuan is small and soft. It has long ears and red eyes. It likes carrots and green vegetables. Every evening, I clean its home and give it fresh water. It can jump very high, and sometimes it jumps onto my bed. My friends all like it, because it is so friendly.\par
Tuantuan is quiet, but it makes my home lively. I love it very much, and it is my best animal friend.\par
\begin{transbox}
现在我家多了一个新成员。它是一只白兔，我们叫它团团。团团又小又软。它有长长的耳朵和红红的眼睛。它喜欢胡萝卜和青菜。每天傍晚，我给它打扫小窝，再换上干净的水。它能跳得很高，有时还会跳到我的床上。我的朋友们也都喜欢它，因为它特别友好。团团很安静，却让我的家变得热闹。我非常爱它，它是我最好的动物朋友。
\end{transbox}

\fanessay{范文七}{我最喜欢的动物}{The Animal I Like Best}{约88词}
My favourite animal is the panda. It is the symbol of China, and it is also the pride of Sichuan.\par
Pandas are black and white. They look like big soft toys, and they have black circles around their eyes. They live in the forests of Sichuan. Every day, they eat bamboo for many hours and sleep for the rest of the day. Baby pandas are very small, but they grow fast.\par
Pandas are cute and peaceful. They need our help, so let's keep their forests green and clean.\par
\begin{transbox}
我最喜欢的动物是熊猫。它是中国的象征，也是四川的骄傲。熊猫是黑白相间的。它们看起来像又大又软的玩具，眼睛四周像画了两个黑圈。它们生活在四川的森林里。每天，它们要吃好几个小时的竹子，剩下的时间都在睡觉。熊猫宝宝非常小，但它们长得很快。熊猫可爱又温和。它们需要我们的帮助，所以让我们守护好它们碧绿、干净的森林吧。
\end{transbox}

\fanessay{范文八}{一次动物园之行}{A Visit to the Zoo}{约81词}
The zoo is my favourite place in Chengdu. I often go there with my family.\par
We usually go in the morning. First, we visit the panda house. The pandas eat bamboo and climb trees. Then we watch the monkeys. They are very naughty, and they make us laugh. My sister likes the giraffes best, because their necks are so long.\par
We go home in the afternoon. I am tired, but I am very happy. I hope to go there again soon.\par
\begin{transbox}
动物园是我在成都最喜欢的地方。我常常和家人一起去那里。我们通常上午去。首先，我们去熊猫馆。熊猫们吃竹子、爬树。然后我们去看猴子。它们非常调皮，逗得我们大笑。我妹妹最喜欢长颈鹿，因为它们的脖子那么长。我们下午回家。我很累，但非常开心。我希望很快能再去那里。
\end{transbox}

\fanessay{范文九}{保护动物，从我做起}{Protecting Animals Starts with Me}{约85词}
Animals are our friends, but many of them are in danger now.\par
People cut down too many trees, so birds lose their homes. Some people kill animals for money. This is very bad. We can help them. We can plant trees in spring. We can keep our parks and rivers clean. We can also tell our friends to be kind to animals.\par
If every one of us does a small thing, animals will have a better life. Let's begin with the small things around us.\par
\begin{transbox}
动物是我们的朋友，但现在许多动物正处在危险之中。人们砍伐太多的树，于是鸟儿失去了家。有些人为了钱而捕杀动物。这非常糟糕。我们可以帮助它们。我们可以在春天种树。我们可以让公园和河流保持干净。我们还可以告诉朋友们要善待动物。如果我们每个人都做一件小事，动物们就会有更好的生活。让我们从身边的小事做起吧。
\end{transbox}

\fanessay{范文十}{我的动物朋友}{My Animal Friend}{约93词}
There is a street cat near my home. It is black and white, and one of its ears is hurt.\par
At first, it runs away when it sees me. Then I put some fish near the wall and wait a few metres away. Now it knows me. Every evening, it waits for me at the same place. It even lets me touch its head.\par
I cannot take it home, because my grandma is afraid of cats. But I can feed it and keep it warm. A friend is a friend, big or small.\par
\begin{transbox}
我家附近有一只流浪猫。它是黑白相间的，一只耳朵受了伤。起初，它一看见我就跑开。后来，我把一些鱼放在墙边，自己站在几米外等着。现在它认识我了。每天傍晚，它都在同一个地方等我。它甚至让我摸摸它的头。我不能把它带回家，因为我奶奶怕猫。但我可以喂它，让它暖和。朋友就是朋友，不分大小。
\end{transbox}


\fanunit{2}{No Rules, No Order}{没有规则就没有秩序}{祈使句；情态动词 must/have to/can；规则表达}

{\centering\small\color{accent}本单元写作任务：My Class Rules（我们的班规）\par}
\vspace{0.6em}

\fansection{〇、本单元写作常用句型}
\begin{fanpatterns}
\fanpat{Every class needs rules, because good rules keep us safe and help us learn better.}{每个班级都需要规则，因为好的规则能保障我们的安全，也能帮助我们更好地学习。}{because 引导原因状语从句，用于开头点明“为什么要有规则”，一句话立住观点。}

\fanpat{Please arrive at school on time and hand in your homework before class.}{请按时到校，并在上课前交作业。}{Please 开头的肯定祈使句，语气客气，适合在介绍规则或提出要求时使用。}

\fanpat{Don't run in the hallway, and don't shout in the reading room.}{不要在走廊上奔跑，也不要在阅览室里大声喊叫。}{Don't 开头的否定祈使句表示“不要做某事”，用 and 并列两条禁令，简洁有力。}

\fanpat{We must wear our school uniforms every Monday, and we have to clean the classroom after school.}{我们每周一必须穿校服，放学后还得打扫教室。}{must 和 have to 都表示“必须”，前者多指规定，后者多指客观要求，常用来陈述校规。}

\fanpat{You can't use your phone in class, but you can use the computer in the computer room.}{你不能在课堂上用手机，但可以在电脑室使用电脑。}{can't 表示“不允许”，can 表示“可以”，一禁一允形成对比，把规则说得清楚明白。}

\fanpat{If everyone follows the rules, our school will be a quiet and happy place to study.}{如果每个人都遵守规则，我们的学校就会成为一个安静、快乐的学习之地。}{if 条件句用一般现在时、主句用一般将来时，用来讲遵守规则的好处，是规则类作文的万能句。}

\fanpat{Not only do rules keep order, but they also teach us to respect each other.}{规则不仅维持秩序，还教会我们彼此尊重。}{Not only 置于句首时主句部分倒装，but... also... 补充递进，适合放在文中或结尾提升档次。}

\fanpat{Let's start with small things and make our school a better place.}{让我们从小事做起，把学校变得更美好。}{Let's 引导的祈使句表示提议和号召，常用于演讲稿或倡议书的结尾。}

\end{fanpatterns}

\fansection{一、满分范文}

\fanessay{范文一}{欢迎来到三班：我们的班规}{Welcome to Class Three: Our Class Rules}{约126词}
Dear Peter, welcome to Class Three! You asked me about our class rules, so let me tell you the most important ones.\par
First, arrive at school before the bell and hand in your homework on time. Second, we must wear our school uniforms every Monday, because we have a flag-raising ceremony that morning. Third, keep the classroom clean and put your rubbish in the bin. Don't run in the hallway, and don't shout in the reading room. You can bring a dictionary to English class, but you can't use your phone during lessons.\par
At first these rules may seem strict, but they help us study well and get on well with each other. Follow them, and you will soon feel at home here. Yours, Li Hua\par
\begin{analysisbox}
1）书信开头用 Dear Peter 称呼，再以 welcome to Class Three 表示欢迎，You asked me about our class rules 交代写信目的，语气友好自然。2）主体用 First、Second、Third 逐条列出班规，must、can、can't 与 Don't 开头的祈使句交替出现，既说清“必须做什么”，也说清“不能做什么”，条理分明。3）结尾先用 At first... but... 承认规则看似严格，再点出好处，最后以 Follow them, and... 收束，把“遵守”与“收获”连在一起，让人乐于接受。
\par\smallskip{\sffamily\bfseries\color{mainblue}亮点句：}{\itshape Follow them, and you will soon feel at home here.}\quad{\small\color{textgray}（译文：遵守这些规则，你很快就会在这里感到像在家里一样自在。）}
\end{analysisbox}

\fanessay{范文二}{一次迟到给我的教训}{A Late Arrival Taught Me a Lesson}{约139词}
Our class has a simple rule: arrive before the bell. Last Tuesday, I broke it.\par
That morning I stopped at the school shop to buy a notebook. The queue was long, and the first bell rang before I could pay. I ran across the playground, but Mr. Li was standing at the door. He asked me to stand at the back of the classroom. My face burned, and I felt everyone was looking at me. I thought the rule was too strict for such a small thing.\par
At lunch, Mr. Li did not scold me. He just said, 'Ten minutes late is still late.' That afternoon I understood something important. A rule is not a cage; it is a promise we make to each other. Now I leave home ten minutes earlier, and I have not been late again.\par
\begin{analysisbox}
1）开头用 Our class has a simple rule 交代班规，紧接着 Last Tuesday, I broke it 形成反差，一句话就把读者带进故事。2）主体按“买本子—排队—铃响—跑到门口—站到教室后面—脸红”的顺序叙述，before I could pay、I ran across the playground 等细节让情节紧凑；I thought the rule was too strict 写出当时的不服气，为下文的转变作铺垫。3）结尾借老师的一句 Ten minutes late is still late 点醒自己，再用 A rule is not a cage; it is a promise we make to each other 升华认识，最后以 I have not been late again 用行动证明改变，真实可信。
\par\smallskip{\sffamily\bfseries\color{mainblue}亮点句：}{\itshape A rule is not a cage; it is a promise we make to each other.}\quad{\small\color{textgray}（译文：规则不是牢笼，而是我们彼此许下的承诺。）}
\end{analysisbox}

\fanessay{范文三}{关于家规，我想和你们聊聊}{Let's Talk About Our Family Rules}{约139词}
Dear Mum and Dad, thank you for taking care of me. I want to talk about our family rules, and I hope you will listen to my ideas.\par
I know rules are important at home, but some of them make me unhappy. For example, I have to finish all my homework before I can play basketball, and I must go to bed at nine every night. Could we change them a little? I would like to play for half an hour first, and then do my homework. On Friday nights, may I go to bed at ten? Also, let's make one rule for the whole family: no phones at the dinner table, so we can talk together.\par
Rules work best when everyone helps to make them. If you agree, let's have a family meeting this Sunday. Yours, Li Hua\par
\begin{analysisbox}
1）书信开头用 Dear Mum and Dad 称呼，thank you for taking care of me 先表达感谢，再提出 I hope you will listen to my ideas，态度礼貌而有主见。2）主体先承认 I know rules are important at home，再用 but some of them make me unhappy 转入建议；have to、must 陈述现有的规定，Could we...? 和 may I...? 用疑问句提出请求，语气委婉，父母更容易接受。3）结尾用 Rules work best when everyone helps to make them 点明道理，并以 If you agree, let's have a family meeting this Sunday 给出具体可行的做法，落点实在。
\par\smallskip{\sffamily\bfseries\color{mainblue}亮点句：}{\itshape Rules work best when everyone helps to make them.}\quad{\small\color{textgray}（译文：当每个人都参与制定规则时，规则才能发挥最好的作用。）}
\end{analysisbox}

\fanessay{范文四}{进图书馆前，请先读这些规则}{Read These Rules Before You Enter the Library}{约116词}
Welcome to our school library! It is a quiet place for reading and studying, so we need some rules to keep it that way. Please read them before you come in.\par
First, keep your voice down. Don't talk loudly or answer phone calls in the reading room. Second, don't bring food or drinks, because they may damage the books. Third, put the books back on the shelf after you read them, and don't move the chairs around. You can borrow two books at a time, and you must return them in two weeks. If you lose a book, please tell the librarian.\par
Follow these rules, and the library will be a comfortable home for every reader.\par
\begin{analysisbox}
1）开头用 Welcome to our school library! 营造告示的语气，so we need some rules to keep it that way 说明立规则的原因，Please read them before you come in 自然引出下文。2）主体用 First、Second、Third 分类说明“保持安静、不带食物饮料、看完归位”等规则，Don't 开头的否定祈使句与 you can、you must 交替使用，既讲清“不能做”，也讲清“可以做什么、必须做什么”，指示明确。3）结尾用“祈使句 + and”的 Follow these rules, and... 收束，把规则与“舒适的阅读之家”联系起来，让人愿意遵守。
\par\smallskip{\sffamily\bfseries\color{mainblue}亮点句：}{\itshape Follow these rules, and the library will be a comfortable home for every reader.}\quad{\small\color{textgray}（译文：遵守这些规则，图书馆就会成为每位读者舒适的家。）}
\end{analysisbox}

\fanessay{范文五}{规则让生活更美好}{Rules Make Life Better}{约134词}
Good morning, everyone! Today I want to talk about rules. Some of you may think rules are boring, but I have found that they make our life better.\par
Think about the traffic lights on Tianfu Avenue. When everybody waits for the green light, the road is safe. Think about a basketball game. Without rules, players push and shout, and no one can enjoy the game. The same is true of our school. When we keep the rules, we have a clean classroom, a quiet library and friendly classmates.\par
So please don't see rules as something that limits us. Rules do not take away our freedom; they keep it safe. They teach us to respect others and help us live together happily. Let's start with small things, and our life will be better and better.\par
\begin{analysisbox}
1）演讲稿开头用 Good morning, everyone! 打招呼，再以 Some of you may think rules are boring, but... 先设对立面再亮出观点，一下子就抓住听众。2）主体用 Think about... 两次举例，把天府大道的红绿灯和篮球比赛摆到听众眼前，再用 when 从句说明规则带来的结果，由交通、比赛推及学校，层层递进。3）结尾用 Rules do not take away our freedom; they keep it safe 一句对比点题，最后以 Let's start with small things... 发出号召，符合演讲的收尾方式。
\par\smallskip{\sffamily\bfseries\color{mainblue}亮点句：}{\itshape Rules do not take away our freedom; they keep it safe.}\quad{\small\color{textgray}（译文：规则并没有夺走我们的自由，而是让自由更安全。）}
\end{analysisbox}

\fansection{二、高分好背好文（适合基础薄弱同学）}

\fanessay{范文六}{我的班规}{My Class Rules}{约85词}
Every class needs rules. Our class has five rules, and I think they are good for us.\par
First, arrive at school on time. Second, wear school uniforms every Monday. Third, keep the classroom clean and tidy. Fourth, don't run in the hallway. Fifth, don't use phones in class. We must also hand in our homework before the first class.\par
At first I thought the rules were too many. Now I know they help us study and keep us safe. I will follow them every day.\par
\begin{transbox}
每个班级都需要规则。我们班有五条规则，我认为它们对我们有好处。第一，按时到校。第二，每周一穿校服。第三，保持教室干净整洁。第四，不要在走廊上奔跑。第五，上课不要用手机。我们还要在第一节课前交作业。起初我觉得规则太多了。现在我知道它们能帮助我们学习，也能保护我们的安全。我每天都会遵守它们。
\end{transbox}

\fanessay{范文七}{我的家规}{My Family Rules}{约89词}
There are some rules in my home, and they are not very strict.\par
I have to get up at seven and make my bed. After school, I must finish my homework before dinner. I can watch TV for half an hour after dinner, but I can't play games on my phone. On Sundays, we clean the house together. My parents also have a rule: no phones at the dinner table.\par
These rules help me use my time well. I think a family with clear rules is a happy family.\par
\begin{transbox}
我家有一些规矩，它们并不很严格。我必须七点起床，并整理好床铺。放学后，我必须在晚饭前完成作业。晚饭后我可以看半小时电视，但不能用手机玩游戏。星期天，我们一起打扫房子。我的父母也有一条规矩：吃饭时不许看手机。这些规矩帮助我很好地利用时间。我认为，一个有明确规矩的家庭是幸福的家庭。
\end{transbox}

\fanessay{范文八}{图书馆规则}{Library Rules}{约88词}
Our school library is a good place to read. Here are some rules for it.\par
Please keep quiet in the reading room. Don't eat or drink there. After you read a book, put it back on the shelf. You can borrow two books at a time. You must return them in two weeks. If you return a book late, you can't borrow books for a week. Don't move the desks or chairs.\par
Follow these rules, and everyone can enjoy the library. Let's keep it clean and quiet together.\par
\begin{transbox}
我们学校的图书馆是一个读书的好地方。下面是我们图书馆的一些规则。请在阅览室保持安静。不要在那里吃东西或喝饮料。读完一本书后，请把它放回书架。你一次可以借两本书。你必须在两周内归还。如果还书晚了，你一周内不能借书。不要搬动桌椅。遵守这些规则，每个人都能享受图书馆。让我们一起保持它的干净和安静。
\end{transbox}

\fanessay{范文九}{我们为什么需要规则}{Why We Need Rules}{约88词}
Some people think rules are boring. But I think rules are important in our life.\par
On the road, rules keep us safe. If everyone runs a red light, accidents will happen. In school, rules keep order. If students talk in class, no one can listen to the teacher. In a game, rules keep it fair. Without them, the strongest always wins, and the game is no fun.\par
Rules are like the lines on a playground. They don't stop us from playing. They make the game fair for everyone.\par
\begin{transbox}
有些人认为规则很无聊。但我认为规则在我们的生活中很重要。在路上，规则保护我们的安全。如果每个人都闯红灯，就会发生事故。在学校，规则维持秩序。如果学生在课堂上讲话，就没有人能听老师讲课。在比赛中，规则保证公平。没有规则，最强的人总是赢，比赛也就没有意思了。规则就像操场上的线。它们并不阻止我们玩耍，而是让比赛对每个人都公平。
\end{transbox}

\fanessay{范文十}{那天，我选择等待}{The Day I Chose to Wait}{约100词}
Last Saturday, my mother and I took the subway to Chunxi Road. The platform was very crowded.\par
When the train came, people pushed to get on. I wanted to push too, because I was afraid of being late. But the rule says, 'Let the passengers get off first.' So I stood still and waited. A woman with a baby in her arms got off slowly. She smiled at me and said, 'Thank you.'\par
We still got on the train and found a seat. That day I learned that following a rule makes others happy, and it makes me happy too.\par
\begin{transbox}
上周六，我和妈妈坐地铁去春熙路。站台上非常拥挤。列车到站时，人们推着往上挤。我也想挤，因为我怕迟到。但规则说：“先下后上。”于是我站着不动，静静等着。一位抱着婴儿的女士慢慢走下车。她对我笑了笑，说：“谢谢你。”我们照样上了车，还找到了座位。那天我明白了，遵守规则能让别人开心，也能让自己开心。
\end{transbox}


\fanunit{3}{Keep Fit}{保持健康}{一般现在时；频率副词；should/can}

{\centering\small\color{accent}本单元写作任务：How I Keep Fit（我如何保持健康）\par}
\vspace{0.6em}

\fansection{〇、本单元写作常用句型}
\begin{fanpatterns}
\fanpat{In my opinion, good health comes from three things: exercise, enough sleep and a cheerful heart.}{在我看来，健康来自三件事：锻炼、充足的睡眠和愉快的心情。}{开头用 In my opinion 亮出观点，再用冒号列出三点，为下文的分类说明搭好框架。}

\fanpat{I go running along the Jinjiang River three times a week.}{我每周沿锦江跑步三次。}{频率副词或频率短语（always, often, twice a week）说明做事的频率，常放在句末或实义动词之前。}

\fanpat{You should get up ten minutes earlier and do some stretching before class.}{你应该早起十分钟，上课前做些拉伸。}{should 加动词原形用来提出建议，语气比 must 委婉，适合给同学或朋友写建议。}

\fanpat{We can keep fit by walking to school instead of taking the bus.}{我们可以用走路上学代替坐公交车来保持健康。}{can 加动词原形表示“可以、能够”，by doing 说明方式，是提建议时最实用的组合。}

\fanpat{Regular exercise not only makes us stronger but also helps us think more clearly.}{经常锻炼不仅让我们更强壮，还能帮助我们思维更清晰。}{not only... but also 连接两个并列成分，用来说明运动的两方面好处，适合放在主体段。}

\fanpat{Two years ago I could hardly run 400 metres, but now I can run 1,000 metres easily.}{两年前我几乎跑不了400米，现在我能轻松跑完1000米。}{用 ago 与 now 对照，配合 could 和 can，写出运动带来的变化，是“运动改变了我”这类作文的核心句。}

\fanpat{Take our class as an example: since we started morning runs, we seldom catch colds.}{以我们班为例：自从开始晨跑，我们很少有人感冒了。}{Take ... as an example 引出具体例子，让倡议或观点更有说服力。}

\fanpat{Let's put down our phones, step out of the classroom and enjoy the sunshine together.}{让我们放下手机，走出教室，一起享受阳光吧。}{Let's 开头的祈使句用于结尾发出号召，简短有力，适合倡议书或演讲稿。}

\end{fanpatterns}

\fansection{一、满分范文}

\fanessay{范文一}{我的锻炼习惯}{My Exercise Habits}{约117词}
Many people think keeping fit means going to the gym, but for me it is simply part of my daily life.\par
I usually get up at six thirty and skip rope for twenty minutes in the yard. Then I walk to school, which takes about fifteen minutes. On Tuesday and Thursday afternoons, I play table tennis with my classmates. I am not a great player, but I never miss a game. On weekends, my father and I often ride bikes along the Tianfu Greenway.\par
These habits take little time, but they keep me energetic all day. I have learned that we do not need a gym to be healthy. We just need to keep moving every day.\par
\begin{analysisbox}
1）开头用 Many people think... but for me... 形成对比，一句话点明“我的锻炼就在日常里”，扣题简洁。2）主体按早晨跳绳、走路上学、下午打球、周末骑行的顺序展开，usually、often 等频率副词和具体时间让习惯真实可信。3）结尾把 take little time 与 keep me energetic 对照，再用 I have learned that 引出感悟，收束自然。
\par\smallskip{\sffamily\bfseries\color{mainblue}亮点句：}{\itshape I have learned that we do not need a gym to be healthy.}\quad{\small\color{textgray}（译文：我明白了，保持健康并不需要健身房。）}
\end{analysisbox}

\fanessay{范文二}{写给同学的健康建议}{Some Advice on Keeping Fit}{约119词}
Dear classmates, I notice that many of us feel tired in class and often catch colds in winter. Here is some advice that may help.\par
First, move more. We should stand up and stretch during the ten-minute break instead of looking at our phones. Second, get enough sleep. Going to bed before ten gives us enough energy for the next day. Third, eat properly. We can bring fruit and milk to school instead of drinking sweet drinks. Finally, stay in a good mood. When we are worried about exams, we should talk to our parents or teachers.\par
Health is the most important thing in our lives. If we follow these small tips, we will study better and smile more.\par
\begin{analysisbox}
1）开头以 Dear classmates 呼告，用 I notice that 句型摆出“上课犯困、冬天常感冒”的现象，自然引出建议。2）主体用 First、Second、Third、Finally 分层，每层都是“观点加做法”，should 与 can 交替出现，建议具体可做。3）结尾先用判断句 Health is the most important thing in our lives 点题，再用 if 条件句展望效果，号召力强。
\par\smallskip{\sffamily\bfseries\color{mainblue}亮点句：}{\itshape If we follow these small tips, we will study better and smile more.}\quad{\small\color{textgray}（译文：如果我们照这些小建议去做，我们会学得更好，笑得更开心。）}
\end{analysisbox}

\fanessay{范文三}{运动会上的八百米}{The 800-Metre Race at the Sports Meeting}{约127词}
Our school sports meeting took place last Friday. I signed up for the 800-metre race, the event I feared most.\par
When the starting gun fired, I ran too fast and soon fell behind. Halfway through, my legs felt as heavy as stones and I wanted to stop. Then I heard my classmates shouting my name along the track. Their voices pulled me forward, and I kept running step by step. In the last hundred metres, I passed two runners and crossed the line in fourth place. I did not win a medal, but I beat my own record by six seconds.\par
Now I understand that a sports meeting is not only about winning. It is about trying your best and knowing that your friends are behind you.\par
\begin{analysisbox}
1）开头用 took place last Friday 交代时间，再用 the event I feared most 制造悬念，读者自然想知道“我”跑得怎样。2）主体按“起跑—落后—听到加油—坚持—冲线”的顺序推进，as heavy as stones、pulled me forward 等表达让画面可感。3）结尾由事及理，先以 not only about winning 递进一层，再用 It is about... 句把运动会写成成长课，升华自然。
\par\smallskip{\sffamily\bfseries\color{mainblue}亮点句：}{\itshape I did not win a medal, but I beat my own record by six seconds.}\quad{\small\color{textgray}（译文：我没有拿到奖牌，但我把自己的纪录提高了六秒。）}
\end{analysisbox}

\fanessay{范文四}{运动改变了我}{Sports Changed Me}{约120词}
Two years ago, I was the quiet boy who sat at the back of the classroom and got out of breath after climbing three floors.\par
Everything changed when my PE teacher asked me to join the running club. At first I could only run 400 metres. I wanted to give up, but he told me to run one more lap each week. Slowly, running became a habit. Now I run for half an hour every evening and can finish 1,000 metres without stopping. I have also made many friends there.\par
Sports have given me not only a stronger body but also a braver heart. I am no longer afraid of difficulties, because I know I can get through them step by step.\par
\begin{analysisbox}
1）开头用 Two years ago 和 who 引导的定语从句描写过去的自己，quiet、out of breath 与后文的改变形成反差。2）主体以 join the running club 为转折，用 At first、Slowly、Now 三个时间词串起变化过程，真实可信。3）结尾用 not only... but also 总结身体与内心的双重收获，braver heart 一语点题。
\par\smallskip{\sffamily\bfseries\color{mainblue}亮点句：}{\itshape Sports have given me not only a stronger body but also a braver heart.}\quad{\small\color{textgray}（译文：运动不仅给了我更强壮的身体，也给了我更勇敢的心。）}
\end{analysisbox}

\fanessay{范文五}{校园健身倡议书}{A Proposal for a Fitter Campus}{约139词}
Dear fellow students, have you noticed that more and more of us spend the break sitting in the classroom? Our eyes are tired and our backs ache, yet we hardly move.\par
I would like to suggest a simple plan. Before the first class, let us run two laps around the playground to wake up our bodies. During the break, we can set up a skipping corner beside the basketball court, and everyone is welcome to join in for five minutes. After school, we should take part in one sport, such as table tennis or badminton. Take our class as an example: since we started morning runs last term, we seldom catch colds.\par
A strong body is the best gift we can give ourselves. Let us put down our phones, step out of the classroom and enjoy the sunshine together.\par
\begin{analysisbox}
1）开头用 have you noticed that 引导的反问句唤起听众注意，tired eyes、aching backs 写出不运动的现状，为倡议作铺垫。2）主体按“课前—课间—课后”三个时间点给出可操作的做法，并用 Take our class as an example 举出本班例子，说服力强。3）结尾用 A strong body is the best gift we can give ourselves 这一富有感染力的句子点明主旨，再以 Let us... 发出号召，语气恳切。
\par\smallskip{\sffamily\bfseries\color{mainblue}亮点句：}{\itshape A strong body is the best gift we can give ourselves.}\quad{\small\color{textgray}（译文：强健的身体是我们能给自己最好的礼物。）}
\end{analysisbox}

\fansection{二、高分好背好文（适合基础薄弱同学）}

\fanessay{范文六}{我如何保持健康}{How I Keep Fit}{约84词}
I am a middle school student in Chengdu. I want to be strong and healthy, so I keep three good habits.\par
First, I exercise every day. I run for twenty minutes in the morning and play basketball on Friday afternoon. Second, I go to bed before ten and get up at six thirty. Third, I eat an apple every day and drink plenty of water.\par
These habits are easy, but they work. I seldom catch a cold now, and I feel happy every day.\par
\begin{transbox}
我是成都的一名中学生。我想变得强壮健康，所以我坚持三个好习惯。第一，我每天都锻炼。早上我跑二十分钟，周五下午打篮球。第二，我十点前睡觉，六点半起床。第三，我每天吃一个苹果，多喝水。这些习惯很简单，却很有效。我现在很少感冒，每天都很快乐。
\end{transbox}

\fanessay{范文七}{我最喜欢的运动}{My Favourite Sport}{约86词}
My favourite sport is table tennis. I started playing it when I was nine.\par
There is a table tennis table in our school playground. I play with my classmates every Tuesday and Thursday after school. The ball is small and fast, so I must watch it carefully. At first I lost almost every game, but now I can win sometimes. My best friend Li Ming often teaches me new skills.\par
Table tennis makes me quick and happy. I hope I can play it for many years.\par
\begin{transbox}
我最喜欢的运动是乒乓球。我九岁开始打乒乓球。我们学校的操场上有一张乒乓球桌。每周二和周四放学后，我和同学们一起打球。球又小又快，所以我必须仔细盯着它。起初我几乎每场都输，但现在我有时也能赢。我最好的朋友李明经常教我新技巧。乒乓球让我反应快、心情好。我希望我能打很多年。
\end{transbox}

\fanessay{范文八}{难忘的一节体育课}{A PE Class to Remember}{约90词}
Last Wednesday, we had a PE class on the playground. The sun was warm and the sky was blue.\par
Our teacher asked us to run 800 metres. I was afraid because I was not good at running. My friend Wang Lei ran beside me and kept saying, “Come on!” On the last lap, I was very tired, but I did not stop. I finished in four minutes.\par
I was the last one in my group, but I was proud of myself. Now I know that finishing is more important than winning.\par
\begin{transbox}
上周三，我们在操场上了一节体育课。阳光暖暖的，天空蓝蓝的。老师让我们跑800米。我很害怕，因为我不擅长跑步。我的朋友王磊在我旁边跑，不停地说：“加油！”最后一圈时，我非常累，但我没有停下来。我用四分钟跑完了。我是我们组最后一个，但我为自己感到骄傲。现在我知道，坚持跑完比赢更重要。
\end{transbox}

\fanessay{范文九}{运动的好处}{Why Sports Are Good for Us}{约83词}
Sports are good for everyone. They help us in many ways.\par
First, sports make our bodies strong. When we run or swim, our hearts and lungs work better. Second, sports make us happy. After a busy day, playing basketball helps me forget my worries. Third, sports teach us to work with others. In a football game, no one can win alone.\par
So let us spend more time doing sports, and start with a ten-minute run today. A healthy body brings a happy life.\par
\begin{transbox}
运动对每个人都有好处。运动在很多方面帮助我们。第一，运动使我们的身体强壮。当我们跑步或游泳时，我们的心肺功能会更好。第二，运动让我们快乐。忙碌一天后，打篮球帮我忘掉烦恼。第三，运动教我们与他人合作。在足球比赛中，没有人能独自取胜。所以让我们多花时间运动，从今天跑十分钟开始吧。健康的身体带来快乐的生活。
\end{transbox}

\fanessay{范文十}{早睡早起}{Early to Bed, Early to Rise}{约86词}
Last term, I often stayed up late to play games. In the morning I could not get up, and I felt sleepy in class.\par
My mother told me to give up the habit. Now I go to bed at ten and get up at six thirty. Before breakfast, I read English for twenty minutes. I also run around the playground twice.\par
Early to bed and early to rise makes me full of energy. I do better in my studies and I am never late for school.\par
\begin{transbox}
上个学期，我常常熬夜打游戏。早上我起不来，上课时觉得很困。妈妈让我改掉这个习惯。现在我十点睡觉，六点半起床。早饭前，我读二十分钟英语。我还绕着操场跑两圈。早睡早起让我精力充沛。我的学习更好了，上学也从不迟到。
\end{transbox}


\fanunit{4}{Eat Well}{吃得健康}{可数与不可数名词；there be；量词}

{\centering\small\color{accent}本单元写作任务：Healthy Eating（健康饮食）\par}
\vspace{0.6em}

\fansection{〇、本单元写作常用句型}
\begin{fanpatterns}
\fanpat{There is some milk in the fridge, and there are some eggs in the basket.}{冰箱里有一些牛奶，篮子里有一些鸡蛋。}{there be 句型表示“某处有某物”，be 动词与最靠近的主语一致；some 多用于肯定句，可修饰不可数名词或复数名词。}

\fanpat{I have a bowl of porridge, an egg and a glass of milk for breakfast.}{我早餐吃一碗粥、一个鸡蛋，喝一杯牛奶。}{不可数名词借助 a bowl of、a glass of、a piece of 等量词表示数量，是介绍一日三餐最常用的写法。}

\fanpat{How much water should we drink every day?}{我们每天应该喝多少水？}{how much 后接不可数名词提问数量，how many 后接可数名词复数，常用于饮食调查或写作开头。}

\fanpat{Eating too much junk food is bad for our health.}{吃太多垃圾食品对我们的健康有害。}{动名词短语作主语时谓语动词用单数，be bad for 说明危害，适合写问题段落。}

\fanpat{We should eat more vegetables and less meat.}{我们应该多吃蔬菜，少吃肉。}{more 和 less 修饰名词表示“更多/更少”，与 should 连用提出饮食建议，简洁有力。}

\fanpat{It is important for us to keep a balanced diet.}{保持均衡饮食对我们很重要。}{It is + 形容词 + for sb. + to do sth. 用 it 作形式主语，是说明“重要性”的高频句式。}

\fanpat{Although fast food tastes good, I only eat it once a week.}{虽然快餐味道不错，但我一周只吃一次。}{although 引导让步状语从句，后面不再用 but；用来写“喜欢但节制”的对比最自然。}

\fanpat{Not only is the food delicious, but it is also good for our health.}{这种食物不仅美味，而且有益健康。}{Not only 置于句首时主句要用部分倒装，多用于推荐食物或结尾升华。}

\end{fanpatterns}

\fansection{一、满分范文}

\fanessay{范文一}{我的健康食谱}{My Healthy Menu}{约122词}
People often ask me why I am always full of energy. My answer is simple: I eat well every day, and I never skip breakfast.\par
For breakfast, I have a bowl of porridge, an egg and a glass of milk. Milk gives me energy for the whole morning. For lunch, I have a bowl of rice with some vegetables and a little meat. There are always two kinds of green vegetables on my plate. For dinner, I have some soup and a piece of fish. I drink about eight glasses of water a day, and I seldom drink cola.\par
Not only does this diet keep me healthy, but it also helps me study better. Good food, I think, is the best medicine.\par
\begin{analysisbox}
1）开头用 People often ask me why 引出话题，再用 My answer is simple 一句点明主题，自然扣住“健康食谱”。2）主体按早餐、午餐、晚餐的顺序展开，a bowl of、a glass of、a piece of 等量词准确修饰不可数名词，There are always two kinds of green vegetables 用 there be 句型写出荤素搭配。3）结尾用 Not only...but also... 的倒装句总结这种饮食的作用，语气自信，收束有力。
\par\smallskip{\sffamily\bfseries\color{mainblue}亮点句：}{\itshape Not only does this diet keep me healthy, but it also helps me study better.}\quad{\small\color{textgray}（译文：这种饮食不仅让我保持健康，还帮助我学习得更好。）}
\end{analysisbox}

\fanessay{范文二}{给朋友的一封饮食建议信}{A Letter to a Friend About Healthy Eating}{约113词}
Dear Tony, I'm sorry to hear that you often feel tired and can't sleep well. As your friend, I'd like to give you some advice on healthy eating.\par
First, never skip breakfast. A glass of milk and an egg are enough to start your day. Second, eat more vegetables and fruit, and eat less junk food. You should drink at least six glasses of water a day instead of cola. Third, try to have dinner before seven o'clock, because eating too late is bad for your stomach.\par
The more healthy food you eat, the stronger you will become. Start with your next meal, and you will feel the change soon. Yours, Li Hua\par
\begin{analysisbox}
1）书信格式完整：Dear Tony 开头，Yours, Li Hua 结尾，I'd like to give you some advice on healthy eating 一句点明写信目的。2）主体用 First、Second、Third 分层给出建议，a glass of、six glasses of 等量词让建议具体可操作，eating too late is bad for your stomach 用动名词作主语说明危害。3）结尾用 the more..., the more... 句型预测改变，鼓励对方从下一餐开始，亲切又有说服力。
\par\smallskip{\sffamily\bfseries\color{mainblue}亮点句：}{\itshape The more healthy food you eat, the stronger you will become.}\quad{\small\color{textgray}（译文：你吃的健康食物越多，你就会变得越强壮。）}
\end{analysisbox}

\fanessay{范文三}{难忘的生日家宴}{An Unforgettable Family Dinner}{约128词}
Last Sunday was my grandmother's seventieth birthday. My whole family got together at her small house in Chengdu.\par
There were twelve people around the big table. Grandma cooked a lot of food for us. There was a bowl of mapo tofu and a plate of twice-cooked pork. In the middle of the table, a hotpot was still boiling. It was so spicy that my little cousin's face turned red and he drank two glasses of water at once. We also had a plate of dumplings. My father kept saying that nothing was better than a family dinner. We ate, talked and laughed for two hours.\par
It was not the food but the laughter that made the dinner unforgettable. I hope we can have such warm dinners more often.\par
\begin{analysisbox}
1）开头用 Last Sunday 和 grandmother's seventieth birthday 交代时间与缘由，一句话把读者带入家宴场景。2）主体围绕餐桌展开，There were twelve people 用 there be 的过去式写人数，There was a bowl of mapo tofu 一句连用 there be 与量词介绍菜品，so... that... 结构写出小表弟被辣得脸红的细节。3）结尾用 It was not...but...that... 强调句点出“难忘”的真正原因，情感真挚。
\par\smallskip{\sffamily\bfseries\color{mainblue}亮点句：}{\itshape It was not the food but the laughter that made the dinner unforgettable.}\quad{\small\color{textgray}（译文：让这顿晚餐难忘的不是食物，而是笑声。）}
\end{analysisbox}

\fanessay{范文四}{校园健康饮食倡议书}{A Proposal for Healthy Eating at School}{约131词}
Dear fellow students, do you know what is in your lunch box? Many of us choose fried chicken and cola because they taste good. However, fried chicken has a lot of fat, and cola has a lot of sugar.\par
Here are three suggestions for you. First, eat breakfast every morning. A bowl of noodles or an egg with a glass of milk is enough. Second, eat more vegetables and fruit, and drink water instead of sweet drinks. Third, try not to eat snacks between meals. There are many healthy choices in our dining hall, such as tomato soup and steamed fish.\par
Not only should we eat well, but we should also remind our friends to do so. Healthy eating today will make a healthier you tomorrow. Let's start with today's lunch!\par
\begin{analysisbox}
1）开头用 Do you know 引导的问句引起听众注意，再指出炸鸡脂肪多、可乐糖分多，问题意识鲜明。2）主体用 Here are three suggestions 总起，First、Second、Third 分层，a bowl of、a glass of 等量词让建议可操作，there are many healthy choices 用 there be 句型给出正面例子。3）结尾用 Not only...but also... 的倒装句发出号召，再用 Healthy eating today will make a healthier you tomorrow 展望未来，Let's start... 收束有力。
\par\smallskip{\sffamily\bfseries\color{mainblue}亮点句：}{\itshape Not only should we eat well, but we should also remind our friends to do so.}\quad{\small\color{textgray}（译文：我们不仅自己要吃好，还应该提醒朋友们也这样做。）}
\end{analysisbox}

\fanessay{范文五}{舌尖上的成都}{The Taste of Chengdu}{约130词}
My hometown is Chengdu, a city famous for its food. If you walk down any street, you can smell something delicious.\par
The most famous dish here is hotpot. There is a lot of chilli and pepper in it, so it tastes hot and spicy. Many people drink a bottle of soy milk with it to cool their mouths. We also have dandan noodles, mapo tofu and dumplings in red oil. A bowl of dandan noodles costs only ten yuan, but it can make you happy for the whole day.\par
There is an old saying in Chengdu that food is the most important thing for people. Our food is not only delicious but also full of love. Welcome to Chengdu, and I will show you the real taste of my hometown.\par
\begin{analysisbox}
1）开头用 a city famous for its food 这一同位语点明成都“美食之都”的身份，If you walk down any street 拉近与读者的距离。2）主体先写火锅，There is a lot of chilli and pepper in it 用 there be 修饰不可数名词，再写担担面、麻婆豆腐等小吃，a bottle of、a bowl of 等量词让介绍具体。3）结尾引用“民以食为天”的老话，用 not only...but also... 把味道与情感联系起来，最后发出邀请，热情自然。
\par\smallskip{\sffamily\bfseries\color{mainblue}亮点句：}{\itshape Our food is not only delicious but also full of love.}\quad{\small\color{textgray}（译文：我们的食物不仅美味，而且充满爱。）}
\end{analysisbox}

\fansection{二、高分好背好文（适合基础薄弱同学）}

\fanessay{范文六}{我最喜欢的食物}{My Favourite Food}{约79词}
My favourite food is dumplings. My mother makes them every weekend, and I always help her.\par
There is some meat and cabbage in the filling. First, my mother mixes the filling. Then we sit around the table and wrap the dumplings together. I can make about twenty in an hour. When they are ready, we dip them in a little vinegar.\par
Dumplings are delicious, and they also mean family to me. That is why I love them so much.\par
\begin{transbox}
我最喜欢的食物是饺子。妈妈每个周末都包饺子，我总是给她帮忙。饺子馅里有肉和白菜。首先，妈妈把馅拌好。然后我们围坐在桌旁一起包饺子。我一小时能包大约二十个。煮好后，我们蘸着一点醋吃。饺子很美味，对我来说还意味着家。这就是我这么喜欢它的原因。
\end{transbox}

\fanessay{范文七}{我的早餐}{My Breakfast}{约81词}
I always have breakfast at home at seven o'clock. My mother says breakfast is the most important meal of the day.\par
I usually have a bowl of porridge, an egg and a glass of milk. Sometimes I eat two pieces of bread with some jam. There is a banana on the table too. I don't like coffee, so I never drink it.\par
After breakfast, I feel full of energy and ready for school. A good breakfast makes my whole morning better.\par
\begin{transbox}
我总是七点在家吃早饭。妈妈说早餐是一天中最重要的一餐。我通常吃一碗粥、一个鸡蛋，喝一杯牛奶。有时我吃两片面包，抹一些果酱。桌上还会有一根香蕉。我不喜欢咖啡，所以从不喝它。早饭后，我精力充沛，做好了上学的准备。一顿好早餐让我的整个上午都更美好。
\end{transbox}

\fanessay{范文八}{健康饮食从我做起}{Healthy Eating Starts with Me}{约86词}
Healthy eating is important for everyone. It helps us grow well and study better.\par
We should eat more vegetables and fruit. There are many kinds of fresh vegetables and fruit in the market. They give us vitamins and keep us strong. We should also drink enough water every day. Junk food is delicious, but it is bad for us, so we should eat only a little.\par
In a word, eat well and live well. Let's start with our next meal and keep it up every day.\par
\begin{transbox}
健康饮食对每个人都很重要。它帮助我们健康成长、更好地学习。我们应该多吃蔬菜和水果。市场上有很多种新鲜蔬菜和水果。它们给我们提供维生素，让我们身体强壮。我们还应该每天喝足够的水。垃圾食品很好吃，但对我们有害，所以我们只能吃一点点。总之，吃得好，生活才美好。让我们从下一餐开始，并且每天都坚持下去。
\end{transbox}

\fanessay{范文九}{一次家庭聚餐}{A Family Dinner}{约85词}
Last Saturday, my family had dinner at a small restaurant near our home.\par
There were six people at the table. We ordered a big fish, a plate of chicken and two bowls of soup. My father also asked for a bottle of juice. The food was fresh and delicious. We talked about my school life and laughed a lot.\par
The best part was being together with my family. I will remember that dinner for a long time. I hope we can go there again soon.\par
\begin{transbox}
上周六，我们一家在家附近的一家小餐馆吃晚饭。餐桌旁有六个人。我们点了一条大鱼、一盘鸡肉和两碗汤。爸爸还要了一瓶果汁。饭菜新鲜又美味。我们聊着我的学校生活，笑得很开心。最美好的部分是和家人聚在一起。这顿晚餐我会记很久。我希望我们很快能再去那里。
\end{transbox}

\fanessay{范文十}{水果与蔬菜}{Fruit and Vegetables}{约87词}
I like fruit and vegetables very much. They are good for my health and they taste nice.\par
Every day I eat an apple and a banana. Sometimes I have a bowl of salad with tomatoes and cucumbers. My mother buys fresh vegetables from the market twice a week. There are carrots, cabbages and green beans in our fridge. She cooks them in different ways.\par
Fruit is sweet, and vegetables are fresh. They make our meals colorful and healthy. Eat more of them, and you will be strong.\par
\begin{transbox}
我非常喜欢水果和蔬菜。它们对我的健康有好处，味道也很好。我每天吃一个苹果和一根香蕉。有时我吃一碗西红柿黄瓜沙拉。妈妈每周去市场买两次新鲜蔬菜。我们冰箱里有胡萝卜、卷心菜和青豆。她用不同的方法烹饪它们。水果甜，蔬菜新鲜。它们让我们的餐桌色彩丰富又健康。多吃水果和蔬菜，你就会变得强壮。
\end{transbox}


\fanunit{5}{Here and Now}{此时此刻}{现在进行时（be doing）；描述正在发生的场景}

{\centering\small\color{accent}本单元写作任务：What Are They Doing?（此刻正在发生的事）\par}
\vspace{0.6em}

\fansection{〇、本单元写作常用句型}
\begin{fanpatterns}
\fanpat{Look! The children are flying kites on the playground.}{看！孩子们正在操场上放风筝。}{Look! 是现在进行时最常用的信号词，一句话就能把读者带进此刻正在发生的场景，适合开头。}

\fanpat{What are they doing? Some are playing games, and others are chatting under the tree.}{他们在做什么？一些人在做游戏，另一些人在树下聊天。}{先用 What are they doing? 设问，再用 some... others... 分述不同人的动作，是描写群体场景的常用组合。}

\fanpat{Right now, my mother is cooking in the kitchen while my father is watering the flowers.}{此时此刻，妈妈正在厨房做饭，而爸爸正在浇花。}{right now 是现在进行时的标志时间状语；while 连接两个同时进行的动作，画面更完整。}

\fanpat{It is raining hard outside, and people are hurrying home with their umbrellas.}{外面正下着大雨，人们打着伞匆匆往家赶。}{先用 It is raining 写天气，再用进行时写人们的反应，是雨中场景描写的固定思路。}

\fanpat{There are many people in the park, and everyone is having a good time.}{公园里有很多人，每个人都玩得很开心。}{there be 句型先交代“有谁”，再用现在进行时写“在做什么”，适合段首总写或段末收束。}

\fanpat{The girl in the red dress is dancing to the music.}{穿红裙子的那个女孩正随着音乐跳舞。}{用 in the red dress 这类介词短语后置定语定位人物，画面更具体，也避免句句以 Some 开头。}

\fanpat{What a lively scene it is!}{这是多么热闹的一幅场景啊！}{What 引导的感叹句用来评价场景，短促有力，常放在段末点出氛围。}

\fanpat{Although everyone is busy, no one feels tired.}{虽然每个人都很忙，但没有人觉得累。}{although 引导让步状语从句，用“忙而不累”的对比收尾，自然升华情感。}

\end{fanpatterns}

\fansection{一、满分范文}

\fanessay{范文一}{公园里的热闹午后}{A Lively Afternoon in the Park}{约118词}
It is four o'clock on Saturday afternoon, and the park near my home is full of life. The sun is shining, and a gentle wind is blowing.\par
Look! Some children are flying kites on the grass, and their kites are dancing in the sky. Under a big tree, an old man is playing chess with his friends, while two boys are watching quietly. Near the lake, a young woman is walking her dog, and the dog is running after a ball. On the small square, some women are dancing to the music. They are all smiling happily.\par
What a lively scene it is! Everyone is busy doing what they love, and the whole park is full of joy.\par
\begin{analysisbox}
1）开头用 It is four o'clock... 交代时间，full of life 总写气氛，The sun is shining, and a gentle wind is blowing 两句先画出天气背景。2）主体用 Look! 领起画面，按草地—树下—湖边—广场的顺序移动镜头，while、and 把同时发生的动作连成一片，人物一多也不显乱。3）结尾用 What 感叹句评价场景，再以 Everyone is busy doing... 收束，由景到情，紧扣“此时此刻”。
\par\smallskip{\sffamily\bfseries\color{mainblue}亮点句：}{\itshape Under a big tree, an old man is playing chess with his friends, while two boys are watching quietly.}\quad{\small\color{textgray}（译文：一棵大树下，一位老人正和朋友们下棋，两个男孩在旁边安静地看着。）}
\end{analysisbox}

\fanessay{范文二}{一次特别的课间}{A Special Break}{约137词}
It is ten past ten in the morning, and the break between classes has just begun. Usually the playground is noisy at this time, but today it is quiet, because it is raining outside.\par
In our classroom, everyone is doing something different. Li Ming and Wang Fang are playing chess at the back of the room, and a small crowd is watching them. Two girls are drawing a big rainbow on the blackboard. By the window, Zhang Lan is reading a storybook. Our English teacher is standing at the door, talking with a boy about his homework. Nobody is running or shouting, but the room is full of quiet laughter and soft voices.\par
When the bell rings, the rain stops and the sun comes out. We rush to the playground, laughing. What a special break this is!\par
\begin{analysisbox}
1）开头先点明课间时间，再用 Usually... but today... 的对比制造悬念：外面下雨，操场反常地安静。2）主体把镜头推进教室，用 everyone is doing something different 总起，再分写下棋、画彩虹、看书、师生交谈四个画面，角度各不相同。3）结尾用 When 从句写出天气转晴，最后以感叹句点题，前后照应“特别”二字。
\par\smallskip{\sffamily\bfseries\color{mainblue}亮点句：}{\itshape Usually the playground is noisy at this time, but today it is quiet, because it is raining outside.}\quad{\small\color{textgray}（译文：平常这个时候操场很吵，但今天很安静，因为外面正在下雨。）}
\end{analysisbox}

\fanessay{范文三}{在表妹的生日派对上}{At My Cousin's Birthday Party}{约129词}
Dear Tony, I am writing to you from my cousin Lily's birthday party. How are you these days? Can you guess what everyone is doing right now?\par
The party is at my aunt's home, and the room is full of color and laughter. My aunt is putting fruit on a big plate, and my uncle is cooking chicken in the kitchen. Some kids are playing games on the floor, and two of them are laughing so hard that they cannot speak. Lily is wearing a pink dress and showing us her new bike. Right now, my grandma is cutting the cake, and we are all singing to Lily.\par
I wish you were here with us. Please write back and tell me what you are doing now. Yours, Chen Jie\par
\begin{analysisbox}
1）开头用书信格式，I am writing to you from... 直接交代“我就在现场”，再问候对方并用设问句 Can you guess...? 引出下文，读者有代入感。2）主体按人物分镜头：姑姑摆水果、叔叔烤鸡、孩子们做游戏、表妹展示新车、奶奶切蛋糕，so... that... 一句写出笑声的程度。3）结尾用 I wish you were here 表达遗憾，并请对方回信说说“你此刻在做什么”，与标题、时态呼应。
\par\smallskip{\sffamily\bfseries\color{mainblue}亮点句：}{\itshape Some kids are playing games on the floor, and two of them are laughing so hard that they cannot speak.}\quad{\small\color{textgray}（译文：一些孩子在地板上做游戏，其中两个笑得说不出话来。）}
\end{analysisbox}

\fanessay{范文四}{雨中的一幕}{A Scene in the Rain}{约130词}
It is half past five, and the rain is coming down hard. School is over, but nobody is leaving. We are all standing at the school gate, watching the water running down the street.\par
Outside, people are hurrying home with their umbrellas. A woman is holding a little girl in her arms, and the girl is laughing at the rain. Two boys are running across the road, holding their schoolbags over their heads. A bus is coming slowly, and the water is splashing under its wheels. At the street corner, an old man is selling roasted sweet potatoes, and steam is rising into the cold air.\par
Then a car stops beside me. It is my father! He is waving at me through the rain. Suddenly the cold afternoon feels warm.\par
\begin{analysisbox}
1）开头用 It is half past five 与 the rain is coming down hard 定下画面基调，nobody is leaving 一句写出大家的无奈，也为后文“被接走”作铺垫。2）主体由近及远：抱孩子的妇女、举着书包跑过的男孩、慢慢进站的公交、卖红薯的老人，动作全部用现在进行时，雨中的城市像一幅流动的画。3）结尾用一个短句 It is my father! 制造惊喜，最后以 cold 与 warm 的对比收束，情景交融。
\par\smallskip{\sffamily\bfseries\color{mainblue}亮点句：}{\itshape A woman is holding a little girl in her arms, and the girl is laughing at the rain.}\quad{\small\color{textgray}（译文：一位妇女怀里抱着一个小女孩，小女孩正对着雨笑。）}
\end{analysisbox}

\fanessay{范文五}{校园艺术节现场报道}{A Live Report from the Art Festival}{约127词}
Good afternoon, everyone. This is Li Hua, reporting live from our school art festival. The festival is going on in the school hall, and the whole building is full of music and laughter.\par
On the stage, some girls are dancing in red dresses. Their dance moves are so wonderful that the audience is cheering loudly. In the corner, some students are painting a long picture of our city, and a boy is playing the guitar softly beside them. At the back of the hall, our head teacher is taking photos, and a few parents are selling handmade cards. Everyone is busy, but everyone is smiling.\par
The festival ends at five, but the happy faces will stay in my mind. Thank you for watching. See you next time!\par
\begin{analysisbox}
1）开头用 This is Li Hua, reporting live from... 的电视直播口吻，一句就确定了“现场报道”的文体，is going on 点明活动正在进行。2）主体按舞台—角落—后排的空间顺序推进，舞蹈、绘画、弹吉他、拍照、义卖各有分工，so... that... 写出观众的反应，画面有动有静。3）结尾用 ends at five 与 will stay in my mind 作时间上的对照，再用一句道别收尾，完整保留直播的格式。
\par\smallskip{\sffamily\bfseries\color{mainblue}亮点句：}{\itshape Their dance moves are so wonderful that the audience is cheering loudly.}\quad{\small\color{textgray}（译文：她们的舞步如此精彩，观众正大声喝彩。）}
\end{analysisbox}

\fansection{二、高分好背好文（适合基础薄弱同学）}

\fanessay{范文六}{此刻的校园}{Our School Right Now}{约94词}
It is three o'clock in the afternoon. Our school is very busy now.\par
On the playground, some boys are playing basketball, and two girls are running around the track. In the library, many students are reading quietly. In my classroom, our teacher is writing on the blackboard, and we are listening carefully. In the music room, a girl is playing the piano, and three students are singing along. Near the school gate, a few parents are waiting for their children.\par
Everyone is busy, but everyone is happy. I love our school at this moment.\par
\begin{transbox}
现在是下午三点。我们学校此刻非常忙碌。操场上，一些男孩正在打篮球，两个女孩正绕着跑道跑步。图书馆里，许多学生正在安静地看书。我们教室里，老师正在黑板上写字，我们正在认真听讲。音乐教室里，一个女孩正在弹钢琴，三个学生正在跟着唱。学校大门附近，几位家长正在等孩子。每个人都很忙，但每个人都很开心。我喜欢此时此刻的学校。
\end{transbox}

\fanessay{范文七}{我们一家正在做的事}{What My Family Is Doing}{约83词}
Now it is seven in the evening. Everyone in my family is at home.\par
My father is watching the news on TV. My mother is washing the dishes in the kitchen. My little sister is drawing a picture at the table. I am doing my homework in my room. My grandparents are drinking tea on the sofa. Our cat is sleeping beside them.\par
The news is still on. The house is quiet, but it feels warm. I like this evening with my family.\par
\begin{transbox}
现在是晚上七点。我们一家人都在家。爸爸正在看电视新闻。妈妈正在厨房洗碗。妹妹正在桌边画画。我正在房间里做作业。爷爷奶奶正在沙发上喝茶。我们的猫正在他们旁边睡觉。新闻还在播着。屋子里很安静，但感觉很温暖。我喜欢和家人在一起的这个夜晚。
\end{transbox}

\fanessay{范文八}{忙碌的早晨}{A Busy Morning}{约84词}
It is six thirty in the morning. Everyone at home is busy.\par
My father is making breakfast in the kitchen. My mother is packing our bags. My sister is looking for her shoes under the bed. I am brushing my teeth in the bathroom. Even our dog is running around the room.\par
Now it is seven, and my father is calling us to hurry up. We are running to the bus stop. The school bus is waiting at the corner. What a busy morning!\par
\begin{transbox}
现在是早上六点半。家里每个人都很忙。爸爸正在厨房做早饭。妈妈正在收拾我们的书包。妹妹正在床下找她的鞋子。我正在卫生间刷牙。连我们的狗都在房间里跑来跑去。现在七点了，爸爸正催我们快点。我们正跑向公交车站。校车正在拐角处等着。多么忙碌的一个早晨啊！
\end{transbox}

\fanessay{范文九}{公园里的人们}{People in the Park}{约85词}
It is Sunday morning, and the park is very busy.\par
Some old people are doing tai chi under the trees. A little boy is playing on the slide, and his mother is watching him. A man is walking his dog along the lake. Two women are sitting on a bench and talking. Near the gate, two boys are flying a kite. A girl is eating ice cream on the grass.\par
The sun is shining, and the birds are singing. Everyone is having a good time.\par
\begin{transbox}
今天是星期天上午，公园里非常热闹。一些老人在树下打太极。一个小男孩正在滑梯上玩，他的妈妈正看着他。一个男人正沿着湖边遛狗。两位女士正坐在长椅上聊天。大门附近，两个男孩正在放风筝。一个女孩正在草地上吃冰激凌。阳光灿烂，鸟儿在歌唱。每个人都玩得很开心。
\end{transbox}

\fanessay{范文十}{我此刻的生活}{My Life at This Moment}{约88词}
It is nine o'clock on Sunday evening. I am sitting at my desk in my room.\par
There is a lot of homework on my desk, but I am looking out of the window. It is raining gently, and the street lights are shining on the wet road. My mother is reading in the living room, and my dog is sleeping beside my chair. A small clock is ticking quietly on the wall.\par
I am a little tired, but I feel peaceful. This is my life at this moment.\par
\begin{transbox}
现在是星期天晚上九点。我正坐在房间里的书桌前。桌上有很多作业，但我正望着窗外。雨正轻轻地下着，路灯照在潮湿的路上。妈妈正在客厅看书，我的狗正在我的椅子旁睡觉。墙上的小钟正轻轻地滴答作响。我有点累，但心里很平静。这就是我此时此刻的生活。
\end{transbox}


\fanunit{6}{Rain or Shine}{无论晴雨}{天气表达；if 条件句；一般将来时}

{\centering\small\color{accent}本单元写作任务：The Weather and My Day（天气与我的日子）\par}
\vspace{0.6em}

\fansection{〇、本单元写作常用句型}
\begin{fanpatterns}
\fanpat{What's the weather like in Chengdu in spring? It's warm and often rainy.}{成都春天的天气怎么样？温暖，而且经常下雨。}{询问天气用 What's the weather like...?，回答用 It's + 天气形容词；这是天气类作文最常用的开头问答句。}

\fanpat{It is sunny and warm today, so I will go to the park with my friends.}{今天阳光明媚、天气温暖，所以我要和朋友们去公园。}{用 It is + 天气形容词 + today 描写当天的天气，再用 so 接当天的活动，一句话写清“天气—活动”的关系。}

\fanpat{If it rains tomorrow, we will stay at home and watch a movie.}{如果明天下雨，我们就待在家里看电影。}{if 引导条件状语从句时遵循“主将从现”：从句用一般现在时，主句用一般将来时，是写计划类作文的核心句型。}

\fanpat{It will be cloudy and cool tomorrow, with a little wind in the afternoon.}{明天多云、凉爽，下午有一点风。}{用 It will be + 天气形容词 预报未来的天气，with 短语补充风力、气温等细节，适合写出行建议。}

\fanpat{Because of the heavy rain, our school trip was put off to next Friday.}{因为大雨，我们的学校郊游推迟到了下周五。}{Because of + 名词短语 说明原因，用来交代天气对计划的影响，比 because 从句更简洁。}

\fanpat{You'd better take an umbrella with you because it may rain in the afternoon.}{你最好带把伞，因为下午可能会下雨。}{You'd better do... 用来给对方出行提醒，because 从句补充理由，语气礼貌又明确。}

\fanpat{When the rain stopped, the sun came out and a rainbow appeared in the sky.}{雨停后，太阳出来了，天空中出现了一道彩虹。}{when 引导时间状语从句，按顺序描写天气的变化，让雨天的叙事更有画面感。}

\fanpat{Rain or shine, I will keep smiling and enjoy every day.}{无论晴雨，我都会保持微笑，享受每一天。}{Rain or shine 意为“无论晴雨”，用作结尾句简洁有力，表达积极的生活态度并呼应单元主题。}

\end{fanpatterns}

\fansection{一、满分范文}

\fanessay{范文一}{雨中的一把伞}{An Umbrella in the Rain}{约136词}
Last Tuesday afternoon, the sky suddenly turned dark when I was walking home from school. Soon the rain came down heavily, and I had to stand at the door of a small shop with other people.\par
I was worried because I had no umbrella and my mother was waiting at home. Just then, a girl in our school uniform opened her umbrella and asked where I lived. She walked with me all the way to my building, though her home was in the other direction. When we said goodbye, the rain was still falling, but my heart felt warm.\par
I never knew her name, but I will always remember that rainy afternoon. If I meet her again, I will say thank you. Rain or shine, a kind heart is the most beautiful thing in the world.\par
\begin{analysisbox}
1）开头用 Last Tuesday afternoon 交代时间，the sky suddenly turned dark 写出天气突变，为下文借伞作铺垫。2）主体按“没带伞—陌生人借伞—绕路送我回家”的顺序叙事，though her home was in the other direction 这一让步细节突出善意，my heart felt warm 用天气与心情的对照收束本段。3）结尾先设想“如果再遇见她”，用 if 条件句表达感谢，再由一次经历上升到对善意的感悟，Rain or shine 呼应单元主题。
\par\smallskip{\sffamily\bfseries\color{mainblue}亮点句：}{\itshape When we said goodbye, the rain was still falling, but my heart felt warm.}\quad{\small\color{textgray}（译文：告别时，雨还在下，但我的心里暖暖的。）}
\end{analysisbox}

\fanessay{范文二}{给来成都的朋友的天气提醒}{A Weather Tip for Your Trip to Chengdu}{约132词}
Dear Mike, I'm glad to hear that you will visit Chengdu next month. Here is some advice about the weather before you start your trip.\par
Spring and autumn are the best seasons here. It is usually warm and sunny, but it can be cool in the morning and at night, so you'd better bring a light jacket. Summer is hot and wet. If you come in July, you will need a hat and an umbrella, because the rain often comes suddenly. Winter is not very cold, but the sky stays grey for weeks, so a warm coat is useful.\par
Chengdu is a city where you can enjoy both sunshine and rain. If you pack the right clothes, you will have a wonderful trip. I'm looking forward to seeing you. Yours, Li Hua\par
\begin{analysisbox}
1）书信格式完整：Dear Mike 开头、Yours, Li Hua 结尾，I'm glad to hear that... 先寒暄，Here is some advice about the weather 一句点明写信目的。2）主体按春夏秋冬的顺序介绍天气，you'd better bring a light jacket 给出穿衣建议，If you come in July, you will need a hat and an umbrella 用“主将从现”的条件句给出具体提醒。3）结尾 If you pack the right clothes, you will have a wonderful trip 再次用条件句收束，Chengdu is a city where... 定语从句让语言更地道。
\par\smallskip{\sffamily\bfseries\color{mainblue}亮点句：}{\itshape If you come in July, you will need a hat and an umbrella, because the rain often comes suddenly.}\quad{\small\color{textgray}（译文：如果你七月来，你需要一顶帽子和一把伞，因为雨常常来得突然。）}
\end{analysisbox}

\fanessay{范文三}{成都的四季天气}{The Weather in Chengdu in Four Seasons}{约134词}
Chengdu is my hometown, and its weather is different in each season. Many people say the city is always cloudy, but that is not the whole story.\par
Spring comes early here. In March, the days get warmer, and flowers bloom along the streets. Summer is hot and rainy. Sometimes the rain lasts for days, and the air feels wet. Autumn is the best season of the year: the sky is clear, the temperature is comfortable, and the leaves turn yellow. In winter it is not very cold. There is little snow, and the sun seldom comes out.\par
Because of the wet weather, people here love tea and hotpot. If you visit Chengdu in autumn, you will understand why life here is so comfortable. The weather may change, but our love for life never does.\par
\begin{analysisbox}
1）开头先承认“多云”的普遍印象，再用 but that is not the whole story 转折，自然引出四季各有特点的主体。2）主体按春、夏、秋、冬的顺序介绍，the temperature is comfortable、the sun seldom comes out 等表达具体自然，冒号后三个并列分句让秋季描写层次清楚。3）结尾用 If you visit Chengdu in autumn, you will understand... 的条件句与 The weather may change, but our love for life never does 的对比句升华主题。
\par\smallskip{\sffamily\bfseries\color{mainblue}亮点句：}{\itshape The weather may change, but our love for life never does.}\quad{\small\color{textgray}（译文：天气会变，但我们对生活的热爱永远不会变。）}
\end{analysisbox}

\fanessay{范文四}{一场雨改变的计划}{When the Rain Changed Our Plan}{约125词}
Last Saturday, my classmates and I planned to climb Qingcheng Mountain. We were going to meet at the school gate at seven in the morning.\par
However, when I woke up, I heard the rain beating on the window. The sky was dark, and the weather report said the rain would last all day. Because of the heavy rain, we had to change our plan. We went to the city library instead. We read books, told jokes and even helped a little boy find his mother. To our surprise, the day was more fun than the climb we had planned.\par
If it is fine next Saturday, we will try the mountain again. But now I know that a rainy day can be a good day too.\par
\begin{analysisbox}
1）开头用 planned to 和 were going to 交代原定计划，meet at seven in the morning 让计划显得具体可信。2）主体用 However 转折，the rain beating on the window 写雨势，Because of the heavy rain 引出计划改变，图书馆里的三件小事说明改了计划也能有收获。3）结尾用 If it is fine next Saturday, we will try the mountain again 的条件句与 But now I know... 的感悟收束，先抑后扬。
\par\smallskip{\sffamily\bfseries\color{mainblue}亮点句：}{\itshape To our surprise, the day was more fun than the climb we had planned.}\quad{\small\color{textgray}（译文：令我们惊讶的是，这一天比我们原本计划的爬山更有趣。）}
\end{analysisbox}

\fanessay{范文五}{天气与心情}{Weather and My Mood}{约133词}
Some people say that the weather decides how we feel. I think it is true in some ways, but not always.\par
On sunny days, I wake up early and want to go outside. The bright sunshine makes me active, and I can finish my homework quickly. On rainy days, I sometimes feel sleepy and lazy, because the grey sky seems to slow everything down. However, I have learned to enjoy rainy days in a different way. I make a cup of hot tea, sit by the window and listen to the rain. Reading a book on such a day is a quiet pleasure.\par
If we choose to be happy, we will find beauty in any weather. Rain or shine, every day is a gift, and it is up to us to enjoy it.\par
\begin{analysisbox}
1）开头用 Some people say... I think... 提出观点，true in some ways, but not always 让态度不绝对化。2）主体以 On sunny days 与 On rainy days 对比两种天气下的状态，再用 However 转入“换个方式过雨天”，hot tea、by the window、listen to the rain 等细节让心情变化真实可感。3）结尾 If we choose to be happy, we will find beauty in any weather 用“主将从现”的条件句点明主旨，Rain or shine 与单元主题呼应。
\par\smallskip{\sffamily\bfseries\color{mainblue}亮点句：}{\itshape If we choose to be happy, we will find beauty in any weather.}\quad{\small\color{textgray}（译文：如果我们选择快乐，我们就会在任何天气里发现美。）}
\end{analysisbox}

\fansection{二、高分好背好文（适合基础薄弱同学）}

\fanessay{范文六}{今天的天气}{The Weather Today}{约89词}
Today was Sunday. The weather was fine. It was sunny and warm.\par
In the morning, I went to the park with my parents. The sky was blue, and there were a few white clouds. Some people were walking, and some were flying kites. We sat on the grass and had a picnic. I ate a sandwich and drank some orange juice.\par
In the afternoon, it became cloudy. It may rain tonight. If it rains, I will stay at home and read a book. I like today's weather very much.\par
\begin{transbox}
今天是星期天。天气很好。阳光明媚，很温暖。上午，我和父母去了公园。天空是蓝色的，还有几朵白云。有些人在散步，有些人在放风筝。我们坐在草地上野餐。我吃了一个三明治，喝了一些橙汁。下午，天变阴了。今晚可能会下雨。如果下雨，我就待在家里看书。我非常喜欢今天的天气。
\end{transbox}

\fanessay{范文七}{我喜欢晴天}{I Love Sunny Days}{约87词}
Different people like different weather. I love sunny days best.\par
On sunny days, the sky is blue and the sun is bright. I feel happy and full of energy. After school, I often play basketball with my friends on the playground. On weekends, my family goes to the park. We take photos and enjoy the sunshine. My little sister likes to run on the grass.\par
Of course, rain is important for plants and for farmers. But if I can choose, I will always choose a sunny day.\par
\begin{transbox}
不同的人喜欢不同的天气。我最喜欢晴天。晴天里，天空是蓝色的，阳光明媚。我感到开心，充满活力。放学后，我常常和朋友们在操场上打篮球。周末，我们一家人去公园。我们拍照，享受阳光。我妹妹喜欢在草地上跑。当然，雨水对植物和农民很重要。但如果我可以选择，我会一直选择晴天。
\end{transbox}

\fanessay{范文八}{一个下雨的下午}{One Rainy Afternoon}{约89词}
It was Tuesday afternoon. When school was over, it was raining hard.\par
I stood at the school gate and looked at the rain. I had no umbrella. My home was about twenty minutes' walk away. I waited for a while, but the rain did not stop. Then my classmate Lucy came. She shared her umbrella with me, and we walked home together. We talked and laughed all the way.\par
When I got home, my shoes were wet, but I was happy. A friend's umbrella made a rainy day warm.\par
\begin{transbox}
那是星期二下午。放学时，雨下得很大。我站在校门口看着雨。我没有带伞。我家大约要走二十分钟。我等了一会儿，但雨没有停。这时，我的同学露西来了。她和我共用一把伞，我们一起走回家。一路上我们说说笑笑。到家时，我的鞋子湿了，但我很开心。朋友的一把伞让雨天变得温暖。
\end{transbox}

\fanessay{范文九}{成都的天气}{The Weather in My City}{约95词}
I live in Chengdu. The weather here is mild and wet, and it changes quickly.\par
In the morning, the sun may be bright, but in the afternoon, dark clouds may cover the sky. Spring and autumn are short and comfortable. Summer is hot and rainy, and winter is grey but not very cold. People here often carry an umbrella, because rain can come at any time.\par
I like this weather, because it makes the city green all year round. If you visit Chengdu, bring a light jacket and an umbrella, and you will be fine.\par
\begin{transbox}
我住在成都。这里的天气温和湿润，而且变化很快。上午可能阳光明媚，到了下午却可能乌云密布。春天和秋天短暂而舒适。夏天炎热多雨，冬天灰蒙蒙的，但不太冷。这里的人常常随身带伞，因为随时可能下雨。我喜欢这样的天气，因为它让这座城市四季常绿。如果你来成都，带上一件薄外套和一把伞，就不会有问题。
\end{transbox}

\fanessay{范文十}{天气与周末计划}{Weather and My Weekend Plan}{约86词}
This weekend, my friends and I have a plan. We want to go to the zoo on Saturday.\par
We are going to meet at nine in the morning. We will take the bus and buy tickets at the gate. I hope the weather will be fine. If it is sunny, we will walk around the zoo and take many photos.\par
But the weather report says it may rain. If it rains, we will go to the science museum instead. We can still have a good time.\par
\begin{transbox}
这个周末，我和朋友们有个计划。我们想星期六去动物园。我们打算早上九点集合。我们会坐公交车，在门口买票。我希望天气好。如果天气晴朗，我们就在动物园里逛逛，拍很多照片。但是天气预报说可能会下雨。如果下雨，我们就改去科学博物馆。我们仍然可以玩得很开心。
\end{transbox}


\fanunit{7}{A Day to Remember}{难忘的一天}{一般过去时；时间顺序与叙事}

{\centering\small\color{accent}本单元写作任务：A Day to Remember（难忘的一天）\par}
\vspace{0.6em}

\fansection{〇、本单元写作常用句型}
\begin{fanpatterns}
\fanpat{Last Sunday was a day I will never forget.}{上个星期天是我永远不会忘记的一天。}{一般过去时开头点题，用时间加感受快速进入叙事，是最常用的记叙文开头。}

\fanpat{First, we cleaned the park; then we planted three young trees.}{首先，我们打扫了公园；然后，我们种了三棵小树。}{first、then、after that、finally 等时间顺序词让叙事层次清楚，动作按先后排列。}

\fanpat{As soon as the bell rang, the whole class rushed out of the classroom.}{铃声一响，全班同学就冲出了教室。}{as soon as 引导时间状语从句，表示“一……就……”，让情节衔接更紧凑。}

\fanpat{While we were waiting for the bus, an old man fell down in front of us.}{我们等公交车的时候，一位老人摔倒在我们面前。}{while 引导的过去进行时从句交代背景，主句用一般过去时写突发动作，画面感强。}

\fanpat{When I heard the good news, my heart jumped with joy.}{听到这个好消息时，我的心高兴得怦怦直跳。}{when 从句加心理描写，写出当时的真实感受，避免只罗列动作。}

\fanpat{Although I was very tired, I felt happier than ever before.}{虽然我很累，但我比以往任何时候都开心。}{although 引导让步状语从句，突出“累却快乐”的反差；注意 although 不能与 but 连用。}

\fanpat{Not only did we help others, but we also learned how to work together.}{我们不仅帮助了别人，还学会了怎样合作。}{Not only 置于句首时主句用部分倒装，but also 补充递进，适合在结尾总结收获。}

\fanpat{From that day on, I have learned to be thankful for what I have.}{从那天起，我学会了对自己拥有的一切心怀感激。}{from that day on 意为“从那天起”，说明事情带来的长久改变，常与现在完成时连用，适合结尾升华。}

\end{fanpatterns}

\fansection{一、满分范文}

\fanessay{范文一}{那个我永远不会忘记的生日}{The Birthday I Will Never Forget}{约137词}
This year my thirteenth birthday fell on a Tuesday. Both of my parents left for work early and said nothing about it. On my way home, I felt sad and thought they had forgotten my birthday.\par
When I opened the door at six, the room was dark. Suddenly the lights came on, and my parents, my grandparents and my best friend all started to sing 'Happy Birthday'. On the table were a big cake and a bowl of dumplings that my grandma had made that afternoon. They told me they had been planning the surprise for two weeks.\par
I made a wish and blew out the candles in one breath. My parents had not forgotten my birthday at all; they just wanted to give me a surprise. That night, I learned that love sometimes hides in silence.\par
\begin{analysisbox}
1）开头没有直接写“难忘”，而是先写父母一早出门、只字未提生日，用 felt sad and thought they had forgotten 制造小小的误会，为后面的惊喜埋下伏笔。2）主体用 When I opened the door at six 把镜头拉进家里，Suddenly the lights came on 一句转折，再用 On the table were... 这一倒装句把蛋糕和饺子推到读者眼前，画面感很强。3）结尾用分号对比“父母并没有忘记，只是想给我惊喜”，最后以 love sometimes hides in silence 点题，把“难忘”落到亲情上，意味深长。
\par\smallskip{\sffamily\bfseries\color{mainblue}亮点句：}{\itshape That night, I learned that love sometimes hides in silence.}\quad{\small\color{textgray}（译文：那天晚上我明白了，爱有时藏在沉默里。）}
\end{analysisbox}

\fanessay{范文二}{青城山上特别的一天}{A Special Day on Qingcheng Mountain}{约137词}
Last April, my father took me to Qingcheng Mountain for a one-day trip. We set off at seven, and the air was cool and fresh. I was sure we would have a wonderful day.\par
It began to rain as soon as we reached the halfway point. The stone steps were wet and slippery, and I fell down twice. I wanted to give up. While I was sitting on a rock, an old couple passed by. The old man smiled and told me that the best view was at the top.\par
With his words in mind, I stood up and went on climbing. At the top, the rain stopped and the clouds lifted. Never had I seen such a green and quiet world. That day taught me that the most beautiful view always comes after the hardest climb.\par
\begin{analysisbox}
1）开头用 Last April 交代时间，took me to Qingcheng Mountain 点明地点和人物，再用 I was sure we would have a wonderful day 写出出发时的好心情，与下文的雨形成反差。2）主体按“上山—遇雨—想放弃—老人鼓励”推进，as soon as 写出雨来得突然，While I was sitting on a rock 用过去进行时交代背景，fell、wanted、passed、smiled、told 等过去式动词让情节连贯。3）结尾用 Never had I seen such a green and quiet world 这一倒装句写出登顶后的震撼，最后以 the most beautiful view always comes after the hardest climb 点明“先苦后甜”的道理，首尾呼应。
\par\smallskip{\sffamily\bfseries\color{mainblue}亮点句：}{\itshape Never had I seen such a green and quiet world.}\quad{\small\color{textgray}（译文：我从未见过如此青翠、宁静的世界。）}
\end{analysisbox}

\fanessay{范文三}{考试失利后的那一天}{The Day I Learned from a Failed Exam}{约138词}
I will never forget the Friday when I got my maths exam paper back. The score was sixty-eight, my lowest ever. I put the paper at the bottom of my bag and walked home slowly, afraid to meet my mother's eyes.\par
That evening, my mother found the paper in my schoolbag. She said nothing about my score, which surprised me. Instead, she cooked me a bowl of noodles and sat beside me. We went over every mistake together, and she wrote down three things to do before the next test.\par
Before bed, I made a simple study plan. From that day on, I did ten problems every evening. Two months later, my score went up to ninety-one. The paper is still in my desk, reminding me that a bad day can be the start of a good change.\par
\begin{analysisbox}
1）开头用 I will never forget the Friday when... 直接点题，The score was sixty-eight, my lowest ever 用具体分数把“失利”写实，afraid to meet my mother's eyes 这一细节写出怕被责备的心理。2）主体写母亲发现试卷后的反应：She said nothing about my score, which surprised me 用非限制性定语从句补充“我”的意外，cooked me a bowl of noodles、sat beside me、went over every mistake 三个动作把“不责备、只陪伴”写得温暖可信。3）结尾用“定计划—每天做题—两个月后提高”的时间链交代转变，最后以 reminding me that a bad day can be the start of a good change 收束，把“难忘”升华为成长。
\par\smallskip{\sffamily\bfseries\color{mainblue}亮点句：}{\itshape The paper is still in my desk, reminding me that a bad day can be the start of a good change.}\quad{\small\color{textgray}（译文：那张试卷还在我的书桌里，它提醒我：糟糕的一天也可以是一个好的改变的开始。）}
\end{analysisbox}

\fanessay{范文四}{帮助一位老人的那一天}{The Day I Helped an Old Man}{约139词}
Last Saturday morning, I was riding my bike to the library when I saw an old man fall down near the market. One handle of his shopping bag broke, and apples rolled all over the road.\par
I stopped at once and ran to him. While I was helping him stand up, a young woman stopped and called his son. The old man's knee was hurt, so I picked up the apples and carried the bag for him. Then I walked him to the bus stop and waited with him.\par
He thanked me again and again and put an apple in my hand. By the time I got to the library, it was almost closing time, but I did not mind at all. On my way home, I felt warm inside, because helping others brings its own kind of happiness.\par
\begin{analysisbox}
1）开头用 I was riding my bike to the library when I saw an old man fall down 这一“过去进行时加 when”结构，把偶遇写得突然而自然。2）主体按“停车—扶人—捡苹果—送站”的动作顺序展开，While I was helping him stand up 交代同时发生的事，picked up、carried、walked、waited with 等动词让帮助的过程具体可感。3）结尾用 By the time... 写出“迟到却不觉遗憾”的反差，最后以 helping others brings its own kind of happiness 点明主题，情理交融。
\par\smallskip{\sffamily\bfseries\color{mainblue}亮点句：}{\itshape On my way home, I felt warm inside, because helping others brings its own kind of happiness.}\quad{\small\color{textgray}（译文：回家的路上，我心里暖暖的，因为帮助别人本身就会带来一种快乐。）}
\end{analysisbox}

\fanessay{范文五}{一家人团聚的那一天}{The Day the Whole Family Came Home}{约139词}
Last Mid-Autumn Festival, my uncle's family drove from Chongqing to Chengdu, so our whole family was finally together. My grandmother got up at five to prepare dinner, and my mother and I put up red lanterns.\par
In the evening, ten of us sat around one big table. Not only did we enjoy my grandmother's dishes, but we also shared the happiest moments of the past year. After dinner, we carried mooncakes out to the balcony. While we were watching the full moon, my grandfather told us stories about his childhood, and even my little cousin listened quietly.\par
Before we went to bed, we took a family photo, which now hangs in our living room. When I look at it, I always remember that warm night. A family does not need to be rich; it only needs to be together.\par
\begin{analysisbox}
1）开头用 Last Mid-Autumn Festival 点明时间，drove from Chongqing to Chengdu、finally together 交代团聚的来之不易，got up at five、put up red lanterns 两个细节写出全家忙碌又欢喜的气氛。2）主体按 In the evening、After dinner 的时间顺序写“吃团圆饭—赏月—听爷爷讲故事”，Not only did we enjoy... 用部分倒装强调物质与情感的双重收获，While we were watching the full moon 用过去进行时烘托温馨气氛。3）结尾用 which now hangs in our living room 这一非限制性定语从句把照片与回忆连起来，最后以 A family does not need to be rich; it only needs to be together 收束，简洁有力。
\par\smallskip{\sffamily\bfseries\color{mainblue}亮点句：}{\itshape Not only did we enjoy my grandmother's dishes, but we also shared the happiest moments of the past year.}\quad{\small\color{textgray}（译文：我们不仅品尝了奶奶做的菜，还分享了过去一年里最开心的时刻。）}
\end{analysisbox}

\fansection{二、高分好背好文（适合基础薄弱同学）}

\fanessay{范文六}{学会骑车的那一天}{The Day I Learned to Ride a Bike}{约93词}
Last summer, my father decided to teach me to ride a bike. I was eleven years old, and I was afraid of falling.\par
On the first day, I fell off three times. My knees hurt, and I wanted to stop. My father did not laugh at me. He held the back of the seat and ran with me again and again. After two hours, I could ride alone.\par
That evening, I rode around the playground and could not stop smiling. I learned that trying one more time can make a hard thing easy.\par
\begin{transbox}
去年夏天，爸爸决定教我骑自行车。那时我十一岁，很怕摔倒。第一天，我摔了三次。膝盖很疼，我想放弃。爸爸没有笑我。他扶着车座，陪我一遍又一遍地跑。两个小时后，我能自己骑了。那天傍晚，我绕着操场骑车，笑个不停。我明白了，再试一次，难事也能变简单。
\end{transbox}

\fanessay{范文七}{我的生日}{My Birthday}{约99词}
Last Sunday was my birthday. I got up early because I was too excited to sleep.\par
In the morning, my mother made a big cake with fruit on it. My father gave me a new football. At noon, three of my friends came to my house, and we played games and sang songs together.\par
In the evening, I made a wish and blew out the candles in one breath. I wished my family and friends would always be happy. Then we ate the cake and took many photos. It was a simple day, but I will always remember it.\par
\begin{transbox}
上周日是我的生日。我起得很早，因为太兴奋了，根本睡不着。上午，妈妈做了一个上面有水果的大蛋糕。爸爸送给我一个新足球。中午，我的三个朋友来到我家，我们一起玩游戏、唱歌。晚上，我许了愿，一口气吹灭了蜡烛。我希望家人和朋友永远快乐。然后我们吃了蛋糕，还拍了很多照片。这是简单的一天，但我会永远记得它。
\end{transbox}

\fanessay{范文八}{一次快乐的周末}{A Happy Weekend}{约93词}
Last weekend was a happy one for me.\par
On Saturday morning, I finished my homework quickly. In the afternoon, my cousin and I went to the park. We flew a kite and ate ice cream. The kite flew higher and higher, and we ran across the grass, laughing.\par
On Sunday, my parents and I visited my grandparents. My grandmother cooked fish, and we talked and laughed. Before we went home, she gave us some fresh vegetables. I was tired on Sunday night, but my heart was full. Weekends with family are the best.\par
\begin{transbox}
上个周末我过得很开心。星期六上午，我很快做完了作业。下午，我和表哥去了公园。我们放风筝、吃冰淇淋。风筝越飞越高，我们在草地上奔跑、欢笑。星期天，我和爸爸妈妈去看望了祖父母。奶奶做了鱼，我们又说又笑。回家前，她给了我们一些新鲜蔬菜。星期天晚上我很累，但心里很满足。和家人在一起的周末最棒。
\end{transbox}

\fanessay{范文九}{糟糕的一天}{A Bad Day}{约104词}
Yesterday was a bad day for me.\par
In the morning, my alarm clock did not ring, so I got up late. I ran to school without breakfast. When I got to class, I found my homework was still at home. My teacher was not happy with me. In the afternoon, it began to rain, and I had no umbrella.\par
But when I opened the door, my mother was cooking hot soup. She smiled and said that a bad day would pass. After dinner, I did my homework and felt better. Now I know that a bad day is only a day, not a life.\par
\begin{transbox}
昨天对我来说是糟糕的一天。早上，我的闹钟没响，所以我起晚了。我没吃早饭就跑去了学校。到了教室，我发现作业还在家里。老师对我不太满意。下午，天开始下雨，而我没有带伞。可是当我打开家门时，妈妈正在煮热汤。她笑着说，糟糕的一天总会过去的。晚饭后，我写完作业，感觉好多了。现在我知道，糟糕的一天只是一天，而不是一生。
\end{transbox}

\fanessay{范文十}{美好的一天}{A Wonderful Day}{约97词}
Last Sunday was a wonderful day. Nothing special happened, but everything went well.\par
I got up at seven and helped my mother water the flowers. Then I went to the bookshop and read a book about space. My friend Li Ming was there, and we talked about the stars.\par
In the evening, my family took a walk by the river. The sky turned orange, and we threw some bread to the fish. On the way home, my father told a funny story, and we all laughed. I did not go anywhere famous, but I felt very happy.\par
\begin{transbox}
上个星期天是美好的一天。没有发生什么特别的事，但一切都顺顺利利。我七点起床，帮妈妈给花浇水。然后我去了书店，读了一本关于太空的书。我的朋友李明也在那里，我们聊了星星。傍晚，我们一家人沿着河边散步。天空变成了橙色，我们往河里扔了些面包喂鱼。回家的路上，爸爸讲了一个好笑的故事，我们都笑了。我没有去任何有名的地方，但我非常开心。
\end{transbox}


\fanunit{8}{Once upon a Time}{很久以前}{一般过去时；故事叙述；once upon a time}

{\centering\small\color{accent}本单元写作任务：A Story I Like（我喜欢的一个故事）\par}
\vspace{0.6em}

\fansection{〇、本单元写作常用句型}
\begin{fanpatterns}
\fanpat{My favourite story is The Tortoise and the Hare, an old fable about a proud hare and a slow tortoise.}{我最喜欢的故事是《龟兔赛跑》，一个关于骄傲的兔子和慢吞吞的乌龟的古老寓言。}{开头点题句式，用同位语补充说明故事内容，一句话交代“写什么”，适合范文第一段。}

\fanpat{Once upon a time, there was a farmer who worked hard in his field every day.}{很久以前，有一位农夫，他每天在田里辛勤劳作。}{once upon a time 是英文故事的经典开头，配合 there be 句型和 who 引导的定语从句介绍人物。}

\fanpat{One day, while the farmer was working in the field, a hare ran into a tree and died.}{一天，农夫正在田里干活时，一只兔子撞到树上死了。}{one day 推进情节，while 引导的过去进行时从句交代背景，是故事叙述最常用的结构之一。}

\fanpat{The hare was so sure of winning that he lay down under a tree and fell asleep.}{兔子太确信自己能赢，于是躺在一棵树下睡着了。}{so... that... 引导结果状语从句，把人物的心理和结果连在一起，让情节发展有因果。}

\fanpat{In the end, the slow tortoise reached the finish line first.}{最后，慢吞吞的乌龟第一个到达了终点。}{in the end 意为“最后”，与 finally 意思相近，用来引出故事的结局，适合放在叙述的收尾句。}

\fanpat{The story tells us that slow and steady wins the race.}{这个故事告诉我们，稳扎稳打才能赢得比赛。}{tell sb. that... 引出寓意，是“故事加道理”类作文最常用的点题句，也可换成 The moral of the story is that...}

\fanpat{From this story, I learned that we should never look down on others.}{从这个故事里，我懂得了我们永远不应该看不起别人。}{From this story, I learned that... 把故事和自己的收获连起来，用于结尾表达启示，注意用一般过去时 learned。}

\fanpat{I'd like to recommend this story to you because it is short but meaningful.}{我想向你推荐这个故事，因为它短小却意味深长。}{I'd like to recommend... to you 是推荐类作文的功能句，because 说明推荐理由，语气礼貌自然。}

\end{fanpatterns}

\fansection{一、满分范文}

\fanessay{范文一}{我最喜欢的故事《守株待兔》}{My Favourite Story: The Farmer and the Hare}{约119词}
My favourite story is The Farmer and the Hare, an old Chinese fable. Once upon a time, there was a farmer who worked hard in his field every day.\par
One day, while he was working, a hare ran into a tree, broke its neck and died. The farmer picked it up happily and took it home for dinner. From that day on, he threw away his hoe and sat by the tree, waiting for another hare. Days passed, but no hare came. In the end, his field was full of weeds.\par
The story tells us that waiting for luck is a waste of time. I like it because it always reminds me to work hard for what I want.\par
\begin{analysisbox}
1）开头用 My favourite story is... 直接点题，再用 Once upon a time 引出故事，符合寓言叙述的经典开头。2）主体按“拾到兔子—放下锄头—田地荒芜”的顺序推进，From that day on、Days passed、In the end 等时间词语让情节清楚，broke its neck、threw away his hoe 等动作描写简练生动。3）结尾用 tells us that 宾语从句点明寓意，再用 reminds me to... 联系自己，主题自然升华。
\par\smallskip{\sffamily\bfseries\color{mainblue}亮点句：}{\itshape The story tells us that waiting for luck is a waste of time.}\quad{\small\color{textgray}（译文：这个故事告诉我们，等待好运只是浪费时间。）}
\end{analysisbox}

\fanessay{范文二}{两粒种子的寓言}{The Fable of Two Seeds}{约125词}
Last week our teacher asked us to write a fable. I made up a story about two seeds. Once upon a time, two seeds lay on the soft soil and talked about their dreams.\par
The first seed said, 'I want to grow into a big tree.' It closed its eyes and went down into the dark earth. The second seed was afraid. 'The soil is dark and cold. I will stay here and enjoy the sun,' it said. So it stayed on the ground. A few days later, a bird came down, picked it up and flew away.\par
The first seed grew into a tall tree, and many birds sang in its branches. My fable tells us that we cannot grow without a little courage.\par
\begin{analysisbox}
1）开头交代写作任务，I made up a story about... 自然引出下文，让“自编寓言”有来处。2）主体用两粒种子的对话和不同选择作对比：一粒勇敢钻进泥土，一粒害怕留在原地，结果一好一坏，对比鲜明。3）结尾先用 grew into a tall tree 写出结果，再用 tells us that 点明寓意，小故事讲出大道理。
\par\smallskip{\sffamily\bfseries\color{mainblue}亮点句：}{\itshape My fable tells us that we cannot grow without a little courage.}\quad{\small\color{textgray}（译文：我的寓言告诉我们，没有一点儿勇气，我们就无法成长。）}
\end{analysisbox}

\fanessay{范文三}{一次难忘的讲故事比赛}{An Unforgettable Storytelling Contest}{约134词}
Last Friday, our school held a storytelling contest in the school hall. I was the fifth speaker, and my hands were cold when I walked onto the stage.\par
I told the story of The Lion and the Mouse. Once upon a time, a lion caught a little mouse and laughed at it. But he let it go. Later, the lion was caught in a net. The little mouse ran to him and bit through the net with its sharp teeth. 'You laughed at me, but I can still help you,' said the mouse. When I finished, the hall was quiet for a second, and then everyone clapped.\par
I did not win first prize, but I learned something important. A small kindness can be a big help, and a quiet story can touch many hearts.\par
\begin{analysisbox}
1）开头用 Last Friday、our school held... 交代时间、地点和事件，my hands were cold 一句写出上场前的紧张，代入感强。2）主体按“上台—讲故事—掌声”的顺序展开，once upon a time 引出故事内容，故事里的对话让叙述更生动。3）结尾用 I did not win first prize, but... 转折写出收获，最后用两个并列句点明主题，立意高于“得奖”。
\par\smallskip{\sffamily\bfseries\color{mainblue}亮点句：}{\itshape A small kindness can be a big help, and a quiet story can touch many hearts.}\quad{\small\color{textgray}（译文：小小的善意也可能帮上大忙，一个安静的故事也能打动许多人的心。）}
\end{analysisbox}

\fanessay{范文四}{愚公移山给我的启示}{What Yu Gong Taught Me}{约127词}
When I was nine, my grandfather told me the story of Yu Gong. Once upon a time, an old man wanted to move two high mountains in front of his house.\par
His neighbours laughed at him. 'You are almost ninety. How can you move a mountain?' Yu Gong answered, 'I will die, but my sons will go on, and their sons after them. The mountains will never grow, but our hands will never stop.' Day after day, he carried stones to the sea. In the end, the gods were moved and took the mountains away.\par
This story gave me a lot of courage. When a maths problem is hard, I think of Yu Gong and try one more time. Big dreams are not built in a day.\par
\begin{analysisbox}
1）开头用 When I was nine, my grandfather told me... 交代故事的来源，亲切自然。2）主体先写邻居的嘲笑，再用愚公的话回应，I will die, but my sons will go on 一句写出“子子孙孙无穷匮也”的精神，In the end 引出结局。3）结尾由故事回到自己的学习生活，用 When a maths problem is hard... 具体说明启示，最后一句点题。
\par\smallskip{\sffamily\bfseries\color{mainblue}亮点句：}{\itshape When a maths problem is hard, I think of Yu Gong and try one more time.}\quad{\small\color{textgray}（译文：遇到解不出的数学题时，我就会想起愚公，再试一次。）}
\end{analysisbox}

\fanessay{范文五}{向你推荐《皇帝的新装》}{A Story I Want You to Read: The Emperor's New Clothes}{约121词}
Dear Peter, I heard that you are looking for an English story to read. I'd like to recommend The Emperor's New Clothes.\par
It is a famous story by Hans Christian Andersen. Once upon a time, an emperor loved new clothes very much. Two cheats told him they could make a magic coat that only clever people could see. The emperor and his officers could see nothing, but no one dared to say so. When the emperor walked through the street, a little child cried out, 'But he has nothing on!'\par
The story is short and easy, and the ending will make you laugh. It also tells us to be honest and brave. I'm sure you will enjoy it. Yours, Li Hua\par
\begin{analysisbox}
1）书信开头用 I heard that you are looking for... 说明写信缘由，I'd like to recommend... 直接点明推荐的书目。2）主体用三句话概括故事：皇帝爱新衣、骗子织“魔衣”、小孩说出真话，情节完整又不啰嗦。3）结尾从“好读”和“有用”两个角度说明推荐理由，I'm sure you will enjoy it 收尾得体。
\par\smallskip{\sffamily\bfseries\color{mainblue}亮点句：}{\itshape It also tells us to be honest and brave.}\quad{\small\color{textgray}（译文：它还告诉我们要诚实、勇敢。）}
\end{analysisbox}

\fansection{二、高分好背好文（适合基础薄弱同学）}

\fanessay{范文六}{龟兔赛跑}{The Tortoise and the Hare}{约94词}
Once upon a time, a hare laughed at a tortoise for being so slow. 'Let's have a race,' said the tortoise.\par
The race began. The hare ran very fast and soon left the tortoise far behind. He thought, 'I have plenty of time. I will have a rest.' So he lay down under a tree and fell asleep. The tortoise kept walking. He did not stop.\par
When the hare woke up, he ran to the finish line, but the tortoise was already there. The story tells us that slow and steady wins the race.\par
\begin{transbox}
很久以前，一只兔子因为乌龟走得太慢而嘲笑它。“我们来比赛吧，”乌龟说。比赛开始了。兔子跑得飞快，不久就把乌龟远远甩在后面。他想：“我有的是时间，我要休息一下。”于是他躺在一棵树下睡着了。乌龟一直走着，没有停下来。兔子醒来时，他跑到终点，可乌龟已经在那儿了。这个故事告诉我们，稳扎稳打才能赢得比赛。
\end{transbox}

\fanessay{范文七}{守株待兔}{Waiting for a Hare by a Tree}{约93词}
Once upon a time, a farmer worked in his field every day. One day, a hare ran into a tree and died.\par
The farmer was very happy. He took the hare home and had a good dinner. Then he said, 'Why work so hard? I can wait for another hare.' He put down his hoe and sat by the tree all day.\par
Days passed, but no hare came. The farmer's field was full of weeds, and his family had nothing to eat. The story tells us that we should not wait for luck.\par
\begin{transbox}
很久以前，一个农夫每天在自己的田里干活。有一天，一只兔子撞到树上死了。农夫非常高兴。他把兔子带回家，美美地吃了一顿晚饭。然后他说：“何必这么辛苦呢？我可以等下一只兔子。”他放下锄头，整天坐在树旁。日子一天天过去，却再也没有兔子来。农夫的田里长满了野草，一家人没有东西吃。这个故事告诉我们，不应该坐等好运降临。
\end{transbox}

\fanessay{范文八}{乌鸦喝水}{The Crow and the Pitcher}{约88词}
My favourite story is The Crow and the Pitcher. It is a short fable about a clever crow.\par
One hot day, a crow was very thirsty. He found a pitcher with a little water in it. His beak could not reach the water. He did not give up. He picked up small stones one by one and dropped them into the pitcher. The water rose higher and higher.\par
At last, the crow drank the water. The story tells us that a clever idea can solve a hard problem.\par
\begin{transbox}
我最喜欢的故事是《乌鸦喝水》。这是一个关于一只聪明乌鸦的短小寓言。炎热的一天，一只乌鸦非常渴。他找到一个罐子，里面有一点儿水。他的嘴够不到水。他没有放弃。他把小石子一颗一颗地捡起来，丢进罐子里。水越升越高。最后，乌鸦喝到了水。这个故事告诉我们，聪明的办法能解决难题。
\end{transbox}

\fanessay{范文九}{灰姑娘}{Cinderella}{约97词}
Once upon a time, there was a kind girl called Cinderella. Her stepmother and two stepsisters made her do all the housework.\par
One day, the prince invited all the girls in the kingdom to a ball. Cinderella could not go, but her fairy godmother helped her. She went to the ball in a beautiful dress and danced with the prince. At twelve o'clock, she ran away and left one glass slipper behind.\par
The prince found the slipper and looked for its owner. It fitted Cinderella. They got married and lived happily together. Kind people deserve good things.\par
\begin{transbox}
很久以前，有一个善良的女孩叫灰姑娘。她的继母和两个继姐妹让她做所有的家务。有一天，王子邀请王国里所有的女孩去参加舞会。灰姑娘不能去，但她的仙女教母帮助了她。她穿着漂亮的裙子去参加舞会，和王子跳舞。十二点时，她跑开了，丢下了一只玻璃鞋。王子找到了那只鞋，寻找它的主人。鞋子正好合灰姑娘的脚。他们结婚了，幸福地生活在一起。善良的人值得拥有美好的事物。
\end{transbox}

\fanessay{范文十}{故事里的道理}{Lessons from Old Stories}{约89词}
I love reading old stories. They are short, but they teach us a lot.\par
The Tortoise and the Hare tells me not to be proud. The Farmer and the Hare tells me not to wait for luck. The Lion and the Mouse tells me that even a little mouse can help a lion.\par
These stories come from long ago, but their lessons are still true today. When I meet a problem, I often ask myself, 'What would the story tell me to do?' Then I know what to do.\par
\begin{transbox}
我喜欢读古老的故事。它们很短，却教给我们很多道理。《龟兔赛跑》告诉我不要骄傲。《守株待兔》告诉我不要等待运气。《狮子与老鼠》告诉我，即使是一只小老鼠也能帮助狮子。这些故事来自很久以前，但其中的道理至今仍然适用。遇到问题时，我常常问自己：“故事会让我怎么做呢？”然后我就知道该怎么做了。
\end{transbox}

\end{document}
